%%%%%%%%%%%%%%%%%%%%%%%%%%%%%%%%%%%%%%%%%%%%%%%%%%%%%%%%%%%%%%%%%%%%%%%%%%%%%%%%%%%%%%%%%%
%%% +++---~~~~ Begin: main document ~~~---+++ %%%%%%%%%%%%%%%%%%%%%%%%%%%%%%%%%%%%%%%%%%%%%%%%%%%%%%
%%%%%%%%%%%%%%%%%%%%%%%%%%%%%%%%%%%%%%%%%%%%%%%%%%%%%%%%%%%%%%%%%%%%%%%%%%%%%%%%%%%%%%%%%%%%%%%%%%%%
%%% Author:         Fabian Roos, fabian@froos.de
%%% Datum:          01.02.2016
%%% Document:       Example Thesis Document
%%% Information:    Compile document with './compile_main_example.sh' or alternatively
%%%                 -> 'pdflatex -shell-escape main_example.tex'
%%%                 -> 'makeglossaries main_example'
%%%                 -> 'bibtex main_example'
%%%                 -> 'pdflatex -shell-escape main_example.tex'
%%%                 -> 'pdflatex -shell-escape main_example.tex'




%%% +++---~~~ Load Header: Settings for the Document ~~~---+++ %%%%%%%%%%%%%%%%%%%%%%%%%%%%%%%%%%%%%
\newcommand{\LANGUAGE}{1}                       % 1 = german, 0 = english
\newcommand{\SHOWDATETIME}{0}                   % 1 = show or 0 = hide the date and time in footer
\newcommand{\TYPE}{0}                           % 1 = master, 0 = bachelor
%%%%%%%%%%%%%%%%%%%%%%%%%%%%%%%%%%%%%%%%%%%%%%%%%%%%%%%%%%%%%%%%%%%%%%%%%%%%%%%%%%%%%%%%%%%%%%%%%%%%
%%% +++---~~~~ start: settings document template: header.tex~~~---+++ %%%%%%%%%%%%%%%%%%%%%%%%%%%%%%
%%%%%%%%%%%%%%%%%%%%%%%%%%%%%%%%%%%%%%%%%%%%%%%%%%%%%%%%%%%%%%%%%%%%%%%%%%%%%%%%%%%%%%%%%%%%%%%%%%%%
%%% Autor:      Fabian Roos, fabian@froos.de
%%% Datum:      01.02.2016 (Nicht Einstellungsänderung, sondern Inhalt!)
%%% Dokument:   Vorlage Masterarbeit (/ Praktikumsbericht) (Vim: kein 'gg=G'!)
%%% Version:    0.1: Erster Entwurf, verwendet für Praktikumsbericht
%%%             0.2: Umstellung auf Paket 'tikzscale'
%%%	            0.3: Restrukturierung der Grafik-Pakete und -Optionen, final für Masterarbeit
%%%             0.4: Einbindung der MWT-Optionen, geometry-Qualmöglichkeit, LaTeX_RS-Version
%%%             0.5: Englische Version für babel, siunitx-Optionen
%%%             0.6: Befehl für Referenzen auf Legenden in PGFPlots
%%%             0.7: Unterstützung für Inkscape-Dateien
%%%             0.8: Flag für Sprache, myunits
%%%             0.9: prepared biber bibliography support
%%%%%%%%%%%%%%%%%%%%%%%%%%%%%%%%%%%%%%%%%%%%%%%%%%%%%%%%%%%%%%%%%%%%%%%%%%%%%%%%%%%%%%%%%%%%%%%%%%%%


%%% +++---~~~ begin: documentclass definitions ~~~---+++ %%%%%%%%%%%%%%%%%%%%%%%%%%%%%%%%%%%%%%%%%%%
% the documentation: http://mirrors.ctan.org/macros/latex/contrib/koma-script/doc/scrguide.pdf
\documentclass[%
% ~~~~~~~~~ Document ~~~~~~~~~~~~~~~~~~~~~~~~~~~~~~~~~~~~~~~~~~~~~~~~~~~~~~~~~~~~~~~~~~~~~~~~~~~~~~~
    % final,    % final document, default option
    % draft,    % show 'overfull hboxes' to see to large elements, no graphics are shown!
% ~~~~~~~~~ font settings ~~~~~~~~~~~~~~~~~~~~~~~~~~~~~~~~~~~~~~~~~~~~~~~~~~~~~~~~~~~~~~~~~~~~~~~~~~
    % fontsize=11pt,    % default value for scrbook, scrreprt, and scrartcl is 11pt, for PHD uses 9pt?
    % headings=normal,  % font size for the headings, which are set large in the scrclass
                        % default should be =big, but =small or =normal are possible as well
% ~~~~~~~~~ Layout ~~~~~~~~~~~~~~~~~~~~~~~~~~~~~~~~~~~~~~~~~~~~~~~~~~~~~~~~~~~~~~~~~~~~~~~~~~~~~~~~~
    paper=a4,	% Einstellung Papiergröße, Standard bei KOMA-Script, s. Handbuch S. 47
    pagesize,	% Schreibt die Dateigröße ins PDF, sollte immer gesetzt sein, s. Handb. S.47
    % landscape,% Landscape-Format
    % oneside,	% Formatierung für einseitiger Druck, Standardoption
    twoside,	% Formatierung für doppelseitiger Druck
    % BCOR=0.78mm,  % Wie viel Platz geht pro Seite für die Bindung verloren? S. 32
    % onecolumn	% Einspaltig, Standard
    % twocolumn,% Zweispaltig
    open=right,	% Wo soll neues Kapitel beginnen? Nur bei scrreprt sinnvoll! S. 89
                % 'any':	Auf jeder neuen Seite
                % 'right':	Auf jeder rechten Seite, fügt ggf. leere Seiten ein!
                %           Ende einer Seite immer auf einer Ebene (\flushbottom wird gesetzt!), S. 29
                %           Dies kann dazu führen, dass es mitunter große Absatzabstände gibt. Daher
                %           ggf. ans Ende einer Seite ein \clearpage http://www.komascript.de/flushbottom
                % 'left':	Auf jeder linken Seite, fügt ggf. leere Seiten ein!
    % parskip=half,     % Absatzabstand /- einzug.
                        % Standard: false, kein Abstand, sondern 1em-Einzug. S.77
    % DIV=calc,   % only for a4 papers default values for the DIV value are defined, if other paper
                % format is chosen, typearea claculates a good value.
                % setting the calculation overwrites the default, use with caution
% ~~~~~~~~ Titelseite ~~~~~~~~~~~~~~~~~~~~~~~~~~~~~~~~~~~~~~~~~~~~~~~~~~~~~~~~~~~~~~~~~~~~~~~~~~~~~~
    % titelpage=false,  % Titelkopf wird gesetzt, Standard bei scrarctl
    % titelpage,        % Titelei in Form von Titelseite, Standard bei scrbook, scrreprt
% ~~~~~~~~ Kopf-/Fuß-Zeile ~~~~~~~~~~~~~~~~~~~~~~~~~~~~~~~~~~~~~~~~~~~~~~~~~~~~~~~~~~~~~~~~~~~~~~~~~
    % headinclude,      % Kopfzeile mit bei der Satzspiegelberechnung beachtet werden? Handb. S. 41
    headsepline,        % Unter Kolumnentitel eine separate Linie setzen, siehe S. 70
    %footsepline,       % Über dem Fuß eine separate Linie setzen, siehe S. 70
% ~~~~~~~~ Verzeichnisse ~~~~~~~~~~~~~~~~~~~~~~~~~~~~~~~~~~~~~~~~~~~~~~~~~~~~~~~~~~~~~~~~~~~~~~~~~~~
    % abstract=true,    % Bei der \begin{abstract}-Umgebung schreibt dies "Zusammenfassung" darüber
    toc=bibliography,   % Literaturverzeichnis ins Inhaltsverzeichnis, S. 73
    % toc=bibliographynumbered,	% Dieses Mal numeriert, S. 73
    % toc=indent,       % Inhaltsverzeichnis hierarchisch, alternativ =graduated, Standard, S. 73
    % toc=flat,         % auch =left, Inhatlsverzeichnis tabellarisch: Gliederungsnummern erste
                        % Spalte, Überschriften zweite, Seitenzahlen dritte Spalte, siehe S. 73
    % toc=index,        % Stichwortverzeichnis in den Index, auch =idx, siehe S. 73
    toc=listof,         % Gleitumgebungen (Abbildungs-, Tabellenverzeichnis) ins Inhaltsverz. S. 72
    % ?                 % weitere Optionen siehe S. 72 und S. 135
    numbers=noenddot,   % Überschritfsnummern ohne Punkt am Ende, alt pointlessnumbers, S. 98
% ~~~~~~~~ Mathematik ~~~~~~~~~~~~~~~~~~~~~~~~~~~~~~~~~~~~~~~~~~~~~~~~~~~~~~~~~~~~~~~~~~~~~~~~~~~~~~
    intlimits,          % place the limits above and below instead of right
    % leqno,            % Formelnummer auf der linken Seite anzeigen, S. 120, amsmath: S. 3
    % fleqn,            % Formeln linksbündig ausrichten, S. 120, amsmath: S. 3
]{scrreprt}             % Komascript-Klassen: scrartcl, scrreprt, scrbook
%%% +++---~~~ Ende: documentclass-Definitionen ~~~---+++  %%%%%%%%%%%%%%%%%%%%%%%%%%%%%%%%%%%%%%%%%%


%%% +++---~~~ Kopfzeile aktivieren ~~~---+++ %%%%%%%%%%%%%%%%%%%%%%%%%%%%%%%%%%%%%%%%%%%%%%%%%%%%%%%
\pagestyle{headings}	% 'plain' ist leere Kopf- und Fußzeile, nur Seitenzahl, Standard, S. 76
                        % 'headings' gibt zusätzlich lebenden Kolumnentitel aus
                        % Besondere Kopfzeilen für u.a. Titelseite, Indexseite: S. 79

% MWT: Einkommentieren: Nur für Datum in Fußzeile? Dann für final auskommentieren?
\usepackage{scrpage2}   % Nur für folgende Definitionen verwenden
\usepackage{scrtime}    % Verschiedene Zeitformate
\pagestyle{scrheadings}
\clearscrheadfoot%      % delete all entries
                        % define header and footer fields
\automark[section]{chapter} % \automark[rechte Seite]{linke Seite} S. 225
\ohead{\headmark}
\ofoot[\pagemark]{\pagemark}
\ifnum\SHOWDATETIME=1
    \ifoot[\tiny\today, \thistime]{\tiny\today, \thistime}
\fi


%%% +++---~~~ Verzeichnis beeinflussen ~~~̣---+++ %%%%%%%%%%%%%%%%%%%%%%%%%%%%%%%%%%%%%%%%%%%%%%%%%%
                        % -1 \part | 0 \chapter | 1 \section | 2 \subsection (Standard)
                        % 3 \subsubsection | 4 paragraph | 5 subparagraph
% \setcounter{tocdepth}{3}      % Bis zu welcher Gliederungsebene im Inhaltsverzeichnis? Siehe S. 70
% \setcounter{secnumdepth}{3}   % Bis zu welcher Gliederungsebene Nummer vor Titel anzeigen? S. 106
                        % Wird ein Kapitel per '\chapter*{Name}' ohne Kapitelnummer erstellt, so
                        % wird es auch nicht im Inhaltsverzeichnis angezeigt.
                        % -> Setze vor die Kapiteldefinition folgende Zeile und das Kapitel taucht
                        %    ToC auf
                        % \addcontentsline{toc}{chapter}{Symbole}
                        % https://tex.stackexchange.com/questions/14379

% MWT: Einkommentieren
% Definition sprachabhängiger Bezeichner, S. 297
\ifcase \LANGUAGE
    \renewcaptionname{english}{\figurename}{Fig.}
    \renewcaptionname{english}{\tablename}{Tab.}
\else
    \renewcaptionname{ngerman}{\figurename}{Bild}
    \renewcaptionname{ngerman}{\tablename}{Tab.}
\fi

% \numberwithin{equation}{section} % (Chapter.Section.Number) statt (Chapter.Number)


%%% +++---~~~ Seitennummerierung beeinflussen ~~~̣---+++ %%%%%%%%%%%%%%%%%%%%%%%%%%%%%%%%%%%%%%%%%%%%
% \pagenumbering{?}     % Setzt die Seitennummerierung auf '1' zurück, wählt Zählweise. Siehe S. 80
                        % arabic (8) / roman (viii) / Roman (VIII) / alph (h) / Alph (H)
                        % Im Dokument an die Stelle setzen, an der die neue Numerierung beginnen
                        % soll
                        % Ggf. muss die Seite erst beendet werden mit '\clearpage' bzw.
                        % '\cleardoublepage'


%%% +++---~~~ Anfang: Pakete laden ~~~---+++ %%%%%%%%%%%%%%%%%%%%%%%%%%%%%%%%%%%%%%%%%%%%%%%%%%%%%%%
% ~~~~~~~~ Standard-Pakete ~~~~~~~~~~~~~~~~~~~~~~~~~~~~~~~~~~~~~~~~~~~~~~~~~~~~~~~~~~~~~~~~~~~~~~~~~
% Für die drei folgenden Pakete siehe auch: http://de.wikipedia.org/wiki/LaTeX#Zeichenkodierung

% <!-- Pakete für compilieren mit pdflatex
\usepackage[utf8]{inputenc}	% Zeichencodierung
\usepackage[T1]{fontenc}	% Schriftart, für Umlaute
% -->

% <!-- Pakete für compilieren mit lualatex
% \usepackage{luainputenc}
% -->

\usepackage{lmodern}        % Hochwertige Fonts: http://tex.stackexchange.com/questions/147194/
\ifcase \LANGUAGE           % load more than one language, active one last
                            % http://tex.stackexchange.com/questions/20987/
    \usepackage[ngerman,english]{babel} % english
\else
    \usepackage[english,ngerman]{babel}	% Für deutsche Begriffe (Inhaltsver., Kapitel) und Worttrennung
                                        % Alte Rechtschreibung: german
                                        % Neue Rechtschreibung: ngerman
                                        % Quelle: http://tex.stackexchange.com/questions/67549/
\fi
% \usepackage[]{csquotes}

% ~~~~~~~~ Grafik-Pakete ~~~~~~~~~~~~~~~~~~~~~~~~~~~~~~~~~~~~~~~~~~~~~~~~~~~~~~~~~~~~~~~~~~~~~~~~~~~
% \usepackage{pdfpages}     % Für das Einbinden von PDF-Dateien als ganze Seite (nicht als Bild)
                            % \includepdf[width=\linewidth]{x.pdf}
\usepackage{graphicx}       % Grafiken: http://en.wikibooks.org/wiki/LaTeX/Importing_Graphics
\usepackage[dvipsnames]{xcolor} % Lade ein definiertes Set an Farben
\usepackage[%               % Beides MWT-Optionen
    hang,
    labelfont = bf,
]{caption}                  % Wird für 'subcaption' benötigt in alten TexLive-Versionen
\usepackage{subcaption}     % Statt veraltete Pakete 'subfigure' / 'subfig'. Siehe Sündenregister
% \usepackage{wrapfig}      % Falls Grafiken von Text umflossen werden sollen, eine Möglichkeit!
                            % Für Grafiken muss Paket graphicx geladen werden
% \usepackage{morefloats}   % Mehr als 18 unbearbeitete floats: http://ctan.org/pkg/morefloats
% \usepackage{flafter}      % MWT: Floats nach Referenz setzen
% \usepackage{placeins}     % MWT: Alle unbearbeiteten Floats anzeigen.
                            % Erlaubt mit \Floatbarrier das Ausgeben aller bisherigen Floats
                            % Achtung: Kann zu teilweise leeren Seiten führen, nur mit Bedacht
                            % verwenden. Besser wäre es mit [!t] oder [!b] zu arbeiten.
\usepackage{adjustbox}      % Paket für die Skalierung von Flussdiagrammen
% \newsavebox{\measurebox}  % verwendet für die Platzierung dreier Bilder nebeneinander

% <!-- Folgende Pakete auch in export.tex
% TikZ
\usepackage{tikz}           % Für Tikz-Grafiken
\usetikzlibrary{%
    external,               % Bibliothek für das externalisieren von TikZ
    calc,                   % Für die Berechnung von Parametern
    shapes.geometric,       % Flussdiagramme
    arrows,                 % Flussdiagramme
    fit,                    % Flussdiagramme: Boxen um Nodes zeichnen
    trees,                  % für 'circular growth' bei den Baumdiagrammen
    % snakes,               % superseded, 'decorations.pathmorphing' verwenden
    decorations.pathmorphing,
}
\usepackage{tkz-euclide}    % Für die Einfache Kennzeichnung von Winkeln
\usetkzobj{all}             % Option für 'tkz-euclide'

% PGF
\usepgflibrary{%
    plotmarks,              % Scatterplots: Kreise statt Linien
}

% PGFPlots (auch in 'export*' benötigt - standalone-Paket-Datei)
\usepackage{pgfplots}       % Paket zum Plotten laden
\usepgfplotslibrary{%
    external,               % Bibliothek für das externalisieren von PGFPlots
    statistics,             % Box Plots
    groupplots,             % Groupplots
    smithchart,             % Smith Charts
    polar,                  % polar plots
    fillbetween,            % fill area with colour
}
\pgfplotsset{compat=newest} % Verwende immer neueste Version und kein Kompatibilitätsmodus
% \pgfplotsset{plot coordinates/math parser=false} % Für matlab2TikZ

% Skalierung von TikZ / PGF
\usepackage{tikzscale}      % Einbinden von TikZ/ PGFPlots zum skalieren per
                            % \includegraphics[width=\linewidth,height=0.9\textheight]{gfx.tikz}
                            % Höhe bei TikZ sinnlos, bei PGFPlots für Achsenverhältnis!
                            % Grafik muss *.tikz sein, Fehler bei *.tex

\usepackage[%
    european,
  ]{circuitikz}             % Für Schaltungen

% Define the Uni Ulm colors
\definecolor{uniblue}{RGB}{125,154,170}
\definecolor{unired}{RGB}{163,38,56}
%
\tikzset{%
    % arrow shape: stealth, latex
    >=stealth,
    %
    %% flow charts
    % -> rectangle blue
    fcrecb/.style={rectangle, rounded corners, align=center, draw=blue, fill=blue!15},
    % -> ellipse red
    fcelr/.style={ellipse, align = center, draw=red, fill=red!20},
    % -> use in document as
    % \tikzset{%
    %   fcblock/.style={fcrecb, minimum width=4cm, minimum height=1.25cm},
    %   fcdecision/.style={fcelr, minimum width=6cm, minimum height=1.25cm},
    % }
}
%
\pgfplotscreateplotcyclelist{roos}{%
  % Usage: Add the following in th axis environment
  % cycle list name = roos,
  % no markers,
  % mark repeat={4},  % Markers only at every nth Datapoint
  % \addplot+
  {blue, solid, mark=x, mark options={solid}},
  {red, densely dashed, mark=o, mark options={solid}},
  {orange, densely dashdotted, mark=triangle, mark options={solid}},
  {black, densely dotted, mark=square, mark options={solid}},
  {brown, loosely dashed, mark=Mercedes star, mark options={solid}},
}
%
\pgfplotscreateplotcyclelist{roosthick}{%
  % Usage: Add the following in th axis environment
  % cycle list name = roos,
  % no markers,
  % mark repeat={4},  % Markers only at every nth Datapoint
  % \addplot+
  {blue, solid, thick, mark=x, mark options={solid}},
  {red, densely dashed, thick, mark=o, mark options={solid}},
  {orange, densely dashdotted, thick, mark=triangle, mark options={solid}},
  {black, densely dotted, thick, mark=square, mark options={solid}},
  {brown, loosely dashed, thick, mark=Mercedes star, mark options={solid}},
}
%
\ifcase \LANGUAGE
\else
% For german documents the decimal seperator is a comma and not a dot. Set this globally
% http://tex.stackexchange.com/a/247423
\pgfplotsset{every linear axis/.append style={%
    /pgf/number format/use comma,
    /pgf/number format/1000 sep={},
  }
}
\fi
%
% Define the parula colormap
\pgfplotsset{% http://tex.stackexchange.com/questions/228197/
    colormap={parula}{%
        rgb = (0.2081,0.1663,0.5292)
        rgb = (0.2116,0.1898,0.5777)
        rgb = (0.2123,0.2138,0.627)
        rgb = (0.2081,0.2386,0.6771)
        rgb = (0.1959,0.2645,0.7279)
        rgb = (0.1707,0.2919,0.7792)
        rgb = (0.1253,0.3242,0.8303)
        rgb = (0.0591,0.3598,0.8683)
        rgb = (0.0117,0.3875,0.882)
        rgb = (0.006,0.4086,0.8828)
        rgb = (0.0165,0.4266,0.8786)
        rgb = (0.0329,0.443,0.872)
        rgb = (0.0498,0.4586,0.8641)
        rgb = (0.0629,0.4737,0.8554)
        rgb = (0.0723,0.4887,0.8467)
        rgb = (0.0779,0.504,0.8384)
        rgb = (0.0793,0.52,0.8312)
        rgb = (0.0749,0.5375,0.8263)
        rgb = (0.0641,0.557,0.824)
        rgb = (0.0488,0.5772,0.8228)
        rgb = (0.0343,0.5966,0.8199)
        rgb = (0.0265,0.6137,0.8135)
        rgb = (0.0239,0.6287,0.8038)
        rgb = (0.0231,0.6418,0.7913)
        rgb = (0.0228,0.6535,0.7768)
        rgb = (0.0267,0.6642,0.7607)
        rgb = (0.0384,0.6743,0.7436)
        rgb = (0.059,0.6838,0.7254)
        rgb = (0.0843,0.6928,0.7062)
        rgb = (0.1133,0.7015,0.6859)
        rgb = (0.1453,0.7098,0.6646)
        rgb = (0.1801,0.7177,0.6424)
        rgb = (0.2178,0.725,0.6193)
        rgb = (0.2586,0.7317,0.5954)
        rgb = (0.3022,0.7376,0.5712)
        rgb = (0.3482,0.7424,0.5473)
        rgb = (0.3953,0.7459,0.5244)
        rgb = (0.442,0.7481,0.5033)
        rgb = (0.4871,0.7491,0.484)
        rgb = (0.53,0.7491,0.4661)
        rgb = (0.5709,0.7485,0.4494)
        rgb = (0.6099,0.7473,0.4337)
        rgb = (0.6473,0.7456,0.4188)
        rgb = (0.6834,0.7435,0.4044)
        rgb = (0.7184,0.7411,0.3905)
        rgb = (0.7525,0.7384,0.3768)
        rgb = (0.7858,0.7356,0.3633)
        rgb = (0.8185,0.7327,0.3498)
        rgb = (0.8507,0.7299,0.336)
        rgb = (0.8824,0.7274,0.3217)
        rgb = (0.9139,0.7258,0.3063)
        rgb = (0.945,0.7261,0.2886)
        rgb = (0.9739,0.7314,0.2666)
        rgb = (0.9938,0.7455,0.2403)
        rgb = (0.999,0.7653,0.2164)
        rgb = (0.9955,0.7861,0.1967)
        rgb = (0.988,0.8066,0.1794)
        rgb = (0.9789,0.8271,0.1633)
        rgb = (0.9697,0.8481,0.1475)
        rgb = (0.9626,0.8705,0.1309)
        rgb = (0.9589,0.8949,0.1132)
        rgb = (0.9598,0.9218,0.0948)
        rgb = (0.9661,0.9514,0.0755)
        rgb = (0.9763,0.9831,0.0538)
    }
}
         % Definitionen für Flussdiagrame (Prefix 'fc'), Plotstyles, uvm.
% Ende der Pakete für export.tex -->

% after a \addplot in a PGFPlots figure a label \label{leg:snr} can be defined and referenced with
% \protect\plotref{leg:snr}
% source: http://tex.stackexchange.com/questions/65494/
% problem: warnings of Multiple defined labels: tex.stackexchange.com/questions/277583/
\newcommand{\plotref}[1]{\tikzexternaldisable\ref{#1}\tikzexternalenable}

% Sollte 'draft' gesetzt werden, so zeige die Boxen in der korrekten Größe an
% Quelle: http://tex.stackexchange.com/questions/47444/
\pgfkeys{/pgf/images/include external/.code=\includegraphics{#1}}

% Optimierung: Aktiviere 'externalize' (Setzen mit --shell-escape)
% \tikzexternalize[%
%     optimize command away=\includepdf, % Eingebundene PDFs nicht externalisieren
%                             % https://tex.stackexchange.com/questions/135504/
%     prefix=tikz/,           % Schreibe externalisierte PDF-Dateien in verzeichnis 'tikz'
% ]

% use lualatex for the externalization process, change inputenc or errors occur
% -> comment in only if you know what you do!
% \tikzset{external/system call={lualatex \tikzexternalcheckshellescape -halt-on-error -interaction=batchmode -jobname "\image" "\texsource"}}

% <!-- Include Inkscape files: Start
% \RequirePackage{scrlfile}                           % prevent the subfig package fromn being loaded
% \PreventPackageFromLoading{subfig}                  % because of known incompatibility with subcaption
%                                                     % use this until the svg package is updated
%                                                     % https://tex.stackexchange.com/questions/291929/
% % \expandafter\def\csname ver@subfig.sty\endcsname{}  % old, not working with Tex Live 2016
%                                                     % http://tex.stackexchange.com/questions/213273/
% \usepackage{svg}                                    % use this package to directly inser svg (!)
%                                                     % files into the document. There is no need for
%                                                     % the pdf_tex approach!
%                                                     % have a look at the README for further
%                                                     % information
% \setsvg{inkscape = inkscape -z -D}                  % with this only the drawing instead of the full
%                                                     % page is included
% \setsvg{svgpath = gfx/}                             % specify the folder where the svg files are
%                                                     % placed
% Include Inkscape files: End -->

\usepackage{xifthen}

% ~~~~~~~~ Mathemathik-Pakete ~~~~~~~~~~~~~~~~~~~~~~~~~~~~~~~~~~~~~~~~~~~~~~~~~~~~~~~~~~~~~~~~~~~~~~
\usepackage{amsmath}        % Für Mathematikumgebungen wie equation, align (& vor =, S. 3)
                            % Formeln referenzieren mit \eqref, siehe S. 10
                            % beinhaltet laut Dokumentation S. 2:
                            % -> amstext: \text innerhalb von Gleichungen
                            % -> amsopn: Neue Formeln mit \DeclareMathOperator
                            % -> amsbsy: Rückwärtskompatibilität
                            % Bei Problemen \usepackage{mathtools}
\usepackage{amssymb}		% Für besondere Zeichen, etwa \dashrightarrow, \mathbb{R}, lädt
                            % -> amsfont: Extra Fonts und Symbole
                            % http://de.wikibooks.org/wiki/LaTeX-Kompendium:_amssymb
\usepackage{trfsigns}       % Für \fourier also o-o
\usepackage{esint}          % special integrals, e.g. \oiin
\ifcase \LANGUAGE
    \usepackage[
        exponent-product={\cdot},       % Produktzeichen bei 1 \cdot 10^-5 ist Standard x (Version 2)
        per-mode=fraction,              % Ohne diese Angabe folgt m s^-1 statt als Bruch
        detect-all,
    ]{siunitx}
\else
    \usepackage[%                       % SIunits ist veraltet, siehe LaTeX2e-Sündenregister S. 15
                                        % Wie soll das Trennzeichen bei Zahlen aussehen?
        % decimalsymbol=comma,	        % Dezimaltrenner (Version 1)
        output-decimal-marker={,},      % Dezimaltrenner (Version 2)
        % expproduct = {\cdot},         % Produktzeichen bei 1 \cdot 10^-5 ist Standard x (Version 1)
        exponent-product={\cdot},       % Produktzeichen bei 1 \cdot 10^-5 ist Standard x (Version 2)
        per-mode=fraction,              % Ohne diese Angabe folgt m s^-1 statt als Bruch
        range-phrase={--},              % Für \SIrange{1}{10}{einheit} statt to ein -
        % range-phrase={ bis },         % Für \SIrange{1}{10}{einheit} statt to ein bis
        %range-units=single,            % Damit nicht nach jedem eine Einheit erscheint
        list-final-separator={ und },   % Für \SIlist{1;2;3}{\einheit}: Letzte Trennung wie?
        list-pair-separator={ und },    % Für \SIlist{1;2}{\einheit}: Trennung in der Mitte wie?
        % list-units=brackets,          % Damit bei SIlist (1;2;3) Einheit erscheint statt 1Einheit, 2Einheit
        detect-all,                     % Damit in strong-Umgebungen auch so dargstellt. Ggf. andere Option?
    ]{siunitx}                          % Beispiel: \SI{1,1}{\metre} kein '{,}' nötig
\fi

\setlength\delimitershortfall{-0.5pt}   % Damit Abstand zwischen mehreren Klammern bei Verwendung
                                        % von \left und \right eingefügt wird. Frank Bögelsack

\allowdisplaybreaks         % Damit in Gleichungen wie align Umbrüche erlaubt sind
                            % http://tex.stackexchange.com/questions/102173/


% ~~~~~~~~ Verschiedenes ~~~~~~~~~~~~~~~~~~~~~~~~~~~~~~~~~~~~~~~~~~~~~~~~~~~~~~~~~~~~~~~~~~~~~~~~~~~
\usepackage{datetime}               % typeset the month name on title page, load after babel!
% \usepackage{cite}                   % Für Mehrfachzitate wie [1-3], sonst nicht benötigt
% \usepackage{csquotes}
% \usepackage{todonotes}              % Bessere (?) Notes, als mit \marginline{...}
% \makeatletter                       % Problem with externalize:
% \renewcommand{\todo}[2][]{\tikzexternaldisable\@todo[#1]{#2}\tikzexternalenable}
% \makeatother                        % http://tex.stackexchange.com/questions/42486/
% \usepackage{blindtext}              % Für Blindtext zum testen, sonst deaktivieren!


% ~~~~~~~~ MWT-Pakete ~~~~~~~~~~~~~~~~~~~~~~~~~~~~~~~~~~~~~~~~~~~~~~~~~~~~~~~~~~~~~~~~~~~~~~~~~~~~~~
\ifcase \LANGUAGE
\else
    \usepackage{bibgerm}            % Falls 'gerunsrt' verwendet werden soll
\fi
% \usepackage{pdflscape}            % Option macht auf Institutsrechner mit externalize Probleme

\usepackage{listings}
\lstset{%
    language=Matlab,                % choose the language of the code
    basicstyle=\scriptsize,         % the size of the fonts that are used for the code
    numbers=none,                   % where to put the line-numbers
    numberstyle=\footnotesize,      % the size of the fonts that are used for the line-numbers
    stepnumber=1,                   % the step between two line-numbers. If it's 1 each line will be numbered
    numbersep=8pt,                  % how far the line-numbers are from the code
    backgroundcolor=\color{white},  % choose the background color. You must add \usepackage{color}
    showspaces=false,               % show spaces adding particular underscores
    showstringspaces=false,         % underline spaces within strings
    showtabs=false,                 % show tabs within strings adding particular underscores
    frame=false,                    % adds a frame around the code
    tabsize=3,                      % sets default tabsize to 2 spaces
    captionpos=b,                   % sets the caption-position to bottom
    breaklines=true,                % sets automatic line breaking
    breakatwhitespace=false,        % sets if automatic breaks should only happen at whitespace
    commentstyle=\color{PineGreen}, % OliveGreen, PineGreen?
    stringstyle=\color{Purple},
    escapeinside={\%*}{*)}          % if you want to add a comment within your code
}

% \usepackage{myunits}                % if siunitx is not used, e.g. 42\Ohm
% \usepackage{mymacros}               % abbrevations

% % Writing floating point numbers in german and math mode (P. Schuh)
% % without additional {}: e.g. $4,56$
% \mathchardef\ocomma="013B
% \mathchardef\pcomma="613B
% \mathcode`,="8000
% {\catcode`,=\active
% \gdef,{\obeyspaces\futurelet\next\smartcomma}}
% \def\smartcomma{\if\space\next\pcomma\else\ocomma\fi}

%% --------------------------------------------------------------------------------------------------
%% Better dealing with floats
%% Alter some LaTeX defaults for better treatment of figures:
%% See p.105 of "TeX Unbound" for suggested values.
%% See pp. 199-200 of Lamport's "LaTeX" book for details.
%%   General parameters, for ALL pages:
%\renewcommand{\topfraction}{0.9}        % max fraction of floats at top
%\renewcommand{\bottomfraction}{0.8}     % max fraction of floats at bottom
%%   Parameters for TEXT pages (not float pages):
%\setcounter{topnumber}{1}
%\setcounter{bottomnumber}{1}
%\setcounter{totalnumber}{4}     % 2 may work better
%\setcounter{dbltopnumber}{2}    % for 2-column pages
%\renewcommand{\dbltopfraction}{0.9}     % fit big float above 2-col. text
%\renewcommand{\textfraction}{0.07}      % allow minimal text w. figs
%%   Parameters for FLOAT pages (not text pages):
%\renewcommand{\floatpagefraction}{0.7}  % require fuller float pages
%% N.B.: floatpagefraction MUST be less than topfraction !!
%\renewcommand{\dblfloatpagefraction}{0.7} % require fuller float pages
%% remember using [htp] or [htpb] for placement
%% --------------------------------------------------------------------------------------------------

% ~~~~~~~~ Layout ~~~~~~~~~~~~~~~~~~~~~~~~~~~~~~~~~~~~~~~~~~~~~~~~~~~~~~~~~~~~~~~~~~~~~~~~~~~~~~~~~~
\usepackage{cmap}           % MWT
\usepackage[activate]{pdfcprot}   % character protruding (MWT)
\usepackage{setspace}       % Falls wirklich notwendig kann hiermit der Zeilenabstand geändert (MWT)
                            % werden. Nicht anders lösen, siehe LaTeX2e-Sündenregister S. 6
\setstretch{1.1}            % skip between lines (MWT)
\usepackage[                % MWT
    textwidth=160mm, textheight=230mm, top=31mm     % A4 Thesis
    % textwidth=111mm, textheight=155mm, top=25mm   % A5 Thesis
]{geometry}

% \usepackage{enumitem}       % u.a. Einzugsänderung und List-Stile ändern
% \usepackage[hyphens]{url}   % Formatieren vorn URLs per \url{http://...}
                            % SEO-URLs haben viele "-" und werden wegen Doppeldeutigkeit nicht
                            % getrennt. Um auch dort zu trennen, die Option [hyphens] übergeben
                            % Achtung: Muss zwingend vor hyperref geladen werden!
                            % http://www.froos.de/blog/latex-url-zu-lang-umbruch-erzwingen-bindestrich/
% \usepackage{float}          % Nötig, um figures per ?H? an Position festzulegen, sonst nicht.
                            % -> Wenn algorithm geladen, dann muss dieses Paket davor geladen
                            %    werden, da algorithm dieses Paket lädt. hyperref braucht ein
                            %    etwaiges float aber davor!
                            % -> tex.stackexchange.com/questions/24586/
\usepackage{hyperref}       % Links im Dokument setzen. Sollte sehr spät geladen werden!
                            % http://en.wikibooks.org/wiki/LaTeX/Hyperlinks
\hypersetup{%
    % hidelinks,            % Keine hevorgehobenen farbigen Links, erst ab Version 6.82a (05.02.2011)
}
% \hypersetup{              % Vorlage für PDF-Eigenschaften, aufnehmen in main.tex
% 	pdftitle={Praktikumsbericht},
% 	pdfsubject={Segmentierung von Laserdaten},
% 	pdfkeywords={Keyword1,Keyword2}
% }

% \usepackage[%
%   style=alphabetic,%
%   sorting=anyt,%
%   sortcites=false,%
%   sortlocale = de,%
%   firstinits=false,%
%   useprefix=false,%
%   uniquename=init,hyperref=auto,%
%   doi=false,isbn=false,url=false,%
%   minnames=3,maxnames=5,%
%   minitems=3,maxitems=99,%
%   minalphanames=3,maxalphanames=3,%
%   maxcitenames=99,%
%   autopunct=false,%
%   backend=biber%
% ]{biblatex}
% \renewcommand*{\labelalphaothers}{}
% \renewcommand\labelnamepunct{\addcomma\space} % https://tex.stackexchange.com/questions/230150/
% \bibliography{literatur}

% ~~~~~~~~ Algorithmen / Pseudocode ~~~~~~~~~~~~~~~~~~~~~~~~~~~~~~~~~~~~~~~~~~~~~~~~~~~~~~~~~~~~~~~~
% Übersicht über mögliche Pakete: http://en.wikibooks.org/wiki/LaTeX/Algorithms
% Welches davon verwenden? tex.stackexchange.com/questions/1375/
% Laut hyperref-Readme erst hyperref und dann algorithm laden
% \usepackage[Pseudocode]{algorithm} % Für eine Algorithmus-Umgebung: Für \label und \caption
% \usepackage{algorithmicx}     % Für die Umgebung des Pseudocodes, definiert Makros
% \usepackage{algpseudocode}    % Formatierung des Pseudocodes, gehört zu algorithmicx
                              % Jeder Pseudocode begint mit \begin{algorithmic}[lines]

% ~~~~~~~~ glossaries and acronyms  ~~~~~~~~~~~~~~~~~~~~~~~~~~~~~~~~~~~~~~~~~~~~~~~~~~~~~~~~~~~~~~~~
% \usepackage[
%     % toc,                        % Im Inhaltsverzeichnis erwähnen
%     % nonumberlist,               % do not show the pages where the entry is used
%                                   % especially do not use if alle entries are printed without \gls{}
%     % nopostdot,                  % no full stop at the end
%     % nogroupskip,                % do not use a gap between the letter groups
%     % xindy,                      % sort with xindy, available with TeX Live
%     % nomain,                     % Falls nicht das Hauptverzeichnis mit 'glo'-Datei verwendet wird
% ]{glossaries}                   % load AFTER hyperref (as mentioned in glossaries manual)
% \setacronymstyle{long-short}
% % \setglossarysection{section}% Festlegen, wie Glossar formatiert wird, Standard ist 'chapter'
% % \newglossary[<log-ext>]{<name>}{<in-ext>}{<out-ext>}{<title>}[<counter>]
% % Quelle: http://www.dickimaw-books.com/latex/thesis/html/makeglossaries.html#sec:newglossary
% % \newglossary[slg]{symblist}{syi}{syg}{Symbolverzeichnis}
% % \newglossary[klg]{konvlist}{kvi}{kvo}{Konventionen}
% \makeglossaries             % Erstellung der externen Dateien für die Indexierung
% % \loadglsentries{symbole}    % Vor '\begin{document}' Symbole laden, aufnehmen in main.tex


%%% +++---~~~ Ende: Pakete laden ~~~---+++ %%%%%%%%%%%%%%%%%%%%%%%%%%%%%%%%%%%%%%%%%%%%%%%%%%%%%%%%%


%%%%%%%%%%%%%%%%%%%%%%%%%%%%%%%%%%%%%%%%%%%%%%%%%%%%%%%%%%%%%%%%%%%%%%%%%%%%%%%%%%%%%%%%%%%%%%%%%%%%
%%% +++---~~~~ Ende: Einstellungen Dokumentenvorlage: header.tex ~~~---+++ %%%%%%%%%%%%%%%%%%%%%%%%%
%%%%%%%%%%%%%%%%%%%%%%%%%%%%%%%%%%%%%%%%%%%%%%%%%%%%%%%%%%%%%%%%%%%%%%%%%%%%%%%%%%%%%%%%%%%%%%%%%%%%
                                  % load the required packages and settings
\newcommand{\MYNAME}{Björn Driemeyer}               % your name
\newcommand{\SUPERVISOR}{Philipp Hügler, M.Sc.} % name of your supervisor
\newcommand{\FIRSTEXAMINER}{Prof.~Dr.-Ing.~Christian Waldschmidt} % your first examiner
\newcommand{\SECONDEXAMINER}{Jun.-Prof.~Dr.-Ing.~Steffen Strehle} % your second examiner
\newcommand{\TITLE}{%                           % this is the title for the title page
  Aufbau und Charakterisierung eines\\ 
  122GHz-Radars 
}

% uncomment and change the following lines for the PDF properties
% \hypersetup{%                                   % define the pdf settings
%   pdfauthor={Fabian Roos},
%   pdftitle={Art der Arbeit},
%   pdfsubject={Längerer Titel der Arbeit},
%   pdfkeywords={keyword1, keyword2}
% }

% \loadglsentries{inhalt/symverz}               % load the glossaries file


%%% +++---~~~ LaTeX: document beginning ~~~---+++ %%%%%%%%%%%%%%%%%%%%%%%%%%%%%%%%%%%%%%%%%%%%%%%%%%
\begin{document}
\pagenumbering{Roman}                           % first page numbers in roman number format


%%% +++---~~~ Titelseite, Sperrverker, Selbstständigkeitserklärung ~~~---+++ %%%%%%%%%%%%%%%%%%%%%%%
\thispagestyle{empty}
\def\unilogohead{%
   \hspace*{-2.2mm}%
   \parbox[b][35mm][t]{\textwidth}{%
      \parbox[t][32mm][t]{\textwidth}{%
      \includegraphics{logos/logo_brief_45_sw}%
      }%
   }%
}%

\begin{center}
\sffamily
\unilogohead
\vspace{10mm}
\Huge\textbf{\TITLE}
\vspace{10mm}

\Large\textbf{%
  \ifcase\LANGUAGE0\ifcase\TYPE0 Bachelor's Thesis\else Master's Thesis\fi
  \else\ifcase\TYPE0 Bachelorarbeit\else Masterarbeit\fi\fi
}
\vspace{20mm}

\huge\textbf{%
\MYNAME%
}
\vfill

{%
\normalsize
\monthname~\the\year%
}

\vfill

\normalsize\textbf{%
\begin{tabular}{rl}
 \ifcase\LANGUAGE0 Supervisor:\else Betreuer:\fi & \SUPERVISOR\\[1ex]
1. \ifcase\LANGUAGE0 Examiner:\else Prüfer:\fi & \FIRSTEXAMINER\\[1ex]
2. \ifcase\LANGUAGE0 Examiner:\else Prüfer:\fi & \SECONDEXAMINER%
\end{tabular}
}
\vspace{15mm}

{\small
\ifcase\LANGUAGE0
Institute of Microwave Engineering\\
Ulm University
\else
Institut für Mikrowellentechnik\\
Universität Ulm
\fi
}

\end{center}


\clearpage
                            % Titelseite laden
\cleardoubleemptypage                           % neue rechte Seite, leere Seite mit Pagestyle empty

% ...oder auch weitere Seiten einbinden

\newcommand{\DECC}{%
Ulm, {\ifcase \LANGUAGE 0 \selectlanguage{english}\else\selectlanguage{ngerman}den \fi\today}
}

%%%%%%%%%%%%%%%%%%%%%%%%%%%%%%%%%%%%%%%%%%%%%%%%%%%%
\newcommand{\DECLARATION}{%
\chapter*{\DECA}
\vspace{2\baselineskip}

\DECB
\vspace{5\baselineskip}

\DECC\hfill\MYNAME
}

\ifcase \LANGUAGE 0
\newcommand{\DECA}{Declaration}
\newcommand{\DECB}{%
I declare that I performed the presented research thesis
%
\begin{center}
\textbf{\large \TITLE}
\end{center}
%
myself and
that I wrote the thesis myself. Any additional sources of information have been
duly identified and referenced. Further, I have regarded the rules of 
scientific work of the Ulm University.
}
\else  %%%%%%%%%%%%%%%%%%%%%%%%%%%%%%%%%%%%%%%%%
\newcommand{\DECA}{Erkl\"{a}rung}
\newcommand{\DECB}{%
Hiermit erkl\"{a}re ich, die vorliegende Arbeit
%
\begin{center}
\textbf{\large \TITLE}
\end{center}
%
selbstst\"{a}ndig verfasst, keine anderen als die angegebenen Quellen
und Hilfsmittel benutzt und die Grunds\"{a}tze und Empfehlungen
\glqq Verantwortung in der Wissenschaft\grqq{} der Universit\"{a}t Ulm
beachtet zu haben.
}
\fi
\DECLARATION
\thispagestyle{empty}

\cleardoubleemptypage                           % neue rechte Seite, leere Seite mit Pagestyle empty


%%% +++---~~~ Define Abstract ~~~---+++ %%%%%%%%%%%%%%%%%%%%%%%%%%%%%%%%%%%%%%%%%%%%%%%%%%%%%%%%%%%
% If you do not know what an abstract is, you don't need it!
% -> if used, activate the abstract in the header for correct formatting
% \begin{abstract}
%   Hier steht die Zusammenfassung des Berichts, realisiert über die \emph{abstract}-Umgebung. Als
%   Überschrift steht dort \emph{Zusammenfassung}.
% \end{abstract}
% \cleardoubleemptypage                           % neue rechte Seite, leere Seite mit Pagestyle empty


%%% +++---~~~ Inhaltsverzeichnis einfügen ~~~---+++ %%%%%%%%%%%%%%%%%%%%%%%%%%%%%%%%%%%%%%%%%%%%%%%
\tableofcontents

\cleardoubleemptypage                           % Damit Seite abgeschlossen ist und die
                                                % Seitennummerierung wieder geändert werden kann


%%% +++---~~~ Dokumentbeginn: Nach Titelseite und Verzeichnisse~~~---+++ %%%%%%%%%%%%%%%%%%%%%%%%%%%
\pagenumbering{arabic}		                    % Ab dem eigentlichen Inhalt
                                                % arabische Numerierung von 1 an

% Inhalte einbinden
%input{inhalt/hilfe}
%\chapter{Radargrundlagen}
In diesem Kapitel werden die für diese Arbeit nötigen Grundlagen hergeleitet und erklärt. 
\section{Radar}
Ein Radar versendet eine elektromagnetische Welle und empfängt, falls ein Objekt in Ausbreitungsrichtung vorhanden ist, die davon reflektierte Welle. Aus Unterschieden zwischen beiden Wellen lassen sich verschiedene Informationen, wie Ort und Geschwindigkeit eines Objekts erkennen. Radar\footnote{\textbf{RA}dio \textbf{D}etection \textbf{A}nd \textbf{R}anging} ist dabei die Bezeichnung für diese Form von Erkennungsverfahren und nutzt dabei drei grundlegende physikalische Gesetze von elektromagnetischen Wellen aus:
\begin{description}
\item Definierte Reflexion an Objekten \\
Wenn ein elektrisch leitendes Objekt getroffen wird, kommt es zu einer Reflexion. Diese neue Welle lässt sich mit einem Empfänger erkennen. 
\item Konstante Geschwindigkeit im Freiraum\footnote{Freiraum bezeichnet den mit Luft gefüllten Bereich zwischen Sender und Zielobjekt} \\
Durch konstante Geschwindigkeit lässt sich die über Messung bis zum Erkennen einer reflektierten elektromagnetischen Welle die Entfernung zu einem Objekt bestimmen. 
\item Geradlinige Ausbreitung \\
Durch spezielle Antennen ist es möglich, Winkel zwischen Sender und Objekt zu bestimmen.
\end{description}
Dabei ist der Begriff Radar nur ein Überbegriff für verschiedene Funktionsweisen die die genannten Gesetzmäßigkeiten nutzen. Anschaulich ist das Radarprinzip in Bild dargestellt.

%%BILD FUER HERLEITUNG

Anschaulich ist das Radarprinzip in Bild dargestellt. Dabei muss zwischen einem monostatischem und einem bistatischem Radar unterschieden werden. In einem monostatischem Radar befinden sich Sender und Empfänger am selben Ort. In einem bistatischem Radar ist der Ort unterschiedlich. \\
Um aus Bild einen Zusammenhang zwischen Sende- und Empfangsleistung zu erhalten wird vorrausgesetzt, dass Sender und Empfänger am gleichen Ort sind, es also ein monostatisches Radar ist. Dabei ist die Sendeleistung $P_{S}$ die Leistung der vom Sender erzeugten Welle, die mit dem Antennengewinn $G_{S}$ abgestrahlt wird. Mit Entfernung von Radar zu Objekt $R$ und dem Rückstrahlfläche $\sigma$ wird eine reflektierte Leistung in entgegengesetzter Richtung, also Richtung Radar, vom Objekt abgestrahlt. Die reflektierte Leistung $P_{R}$ ergibt sich demnach zu
\begin{align}
P_{R} = \frac{P_{S} \cdot G_{S} \cdot \sigma}{4\pi \cdot R^2}.
\end{align} Die Empfangsleistung ergibt sich über die reflektierte Leistung und die Antennenwirkfläche $A$ wobei für $A{W}$ gilt mit dem Antennengewinn der Empfangsantenne $G_{e}$\footnote{In diesem Fall ist die Empfangsantenne und die Sendeantenne identisch, allerdings lässt sich die Herleitung verallgemeinern, wofür diese Unterscheidung notwendig ist}. 
\begin{align}
A = \frac{\lambda^2 \cdot G_{e}}{4\pi}.
\end{align}
Dadurch ergibt sich die über die Empfangsleistungsdichte $S_{E}$, wobei 
\begin{align}
S_{E} = \frac{P_{R}}{4\pi \cdot R^2}.
\end{align}
ist, die 
Empfangsleistung $P_{E}$ mit
\begin{align}
P_{E} = S_{E}\cdot A_{W}.
\end{align}
und mit obigen Gleichungen
\begin{align}
P_{E} = \frac{P_{S} \cdot G_{S} \cdot G_{E} \cdot \sigma \cdot \lambda^2}{(4\pi)^3\cdot R^4}
\end{align}
die Radargleichung\footnote{Einzige Einschränkung ist hier, dass Sender und Empfänger die gleiche Entfernung vom Ziel haben müssen. Mit obigen Einschränkungen ist zusätzlich $G_{S} = G_{E}$ wodurch sich die Radargleichung weiter vereinfacht.}.
Durch Umstellung nach $R$ unter Berücksichtigung der minimale Empfangsleistung $P_{E,min}$\footnote{$P_{E,min}$ ist dabei proportional zur Empfängerbandbreite
$B$, der Empfängerrauschzahl $F$ und dem minimalen Signal-zu-Rausch-Verhältnis $SNR_{min}$.} ergibt sich eine maximale Radarreichweite $R_{max}$ mit
\begin{align}
R_{max} = \sqrt[4]{\frac{P_{S} \cdot G_{E}\cdot G_{S} \cdot \lambda^2 \cdot \sigma}{P_{E,min} \cdot (4\pi)^3}}.
\end{align}
Das Verhältnis von ausgesendeter Leistung und empfangener Leistung bezeichnet die Freiraumdämpfung $F$ oder FSPL\footnote{\textbf{F}ree\textbf{S}pace\textbf{P}ath\textbf{L}oss}. Dabei ist
\begin{align}
F = \frac{P_{E}}{P_{S}} = \left(\frac{4\pi \cdot R}{\lambda}\right)^2.
\end{align} 
und mit Lichtgeschwindigkeit c sowie Frequenz f
\begin{align}
c = \lambda \cdot f.
\label{eq:fzulambda}
\end{align}
die Freiraumdämpfung
\begin{align}
F(f,r) = \left(\frac{4\pi \cdot R \cdot f}{c}\right)^2
\end{align}
eine Funktion in Abhängigkeit der Frequenz, und dem Abstand $R$. Diese Frequenz bleibt konstant. Häufig wird $F(f)$ in dB und mit der englischen Bezeichnung bezeichnet, demnach ist
\begin{align}
\text{FPSL(dB)} &= 10\log_{10}\left(\frac{4\pi \cdot R \cdot f}{c}\right)^2 \notag \\
&= 20 \log_{10}\left( R \right) + 20 \log_{10}\left( f \right) - 147,55 .
\end{align}

\section{Radartypen}
Radarsysteme lassen sich in Impuls- und Dauerstrichradarsysteme unterteilen. Das Impulsradar ist hier nicht beschrieben, da es für die weitere Arbeit keine  Verwendung hat. Als Vertreter der Dauerstrichradare wird besonders auf das FMCW Radar eingegangen.

\subsection{Zielentfernung}
Das Radar sendet eine Signal in den Freiraum aus und misst die Zeit, bis das reflektierte Signal detektiert wird. Dieser Zeitunterschied $\Delta t$ gibt mit Hilfe der konstanten Lichtgeschwindigkeit $c$ den Abstand $R$ von Radar zu Objekt über
\begin{align}
R = \frac{c \cdot \Delta t}{2} 
\end{align} 
an. Bezieht man zusätzlich eine konstante Radialgeschwindigkeit des Ziels mit ein, gilt für den Abstand 
\begin{equation}
R(t) = R_{0} - v_{r} \cdot \left( t-t_{0}\right).
\end{equation}
Für den Laufzeitunterschied gilt demnach nach \cite[Gl. 3.1.11]{HuderRadar}
\begin{equation}
\Delta t = \frac{2R}{c_{0}}
		 = \frac{2}{c_{0}}\left(R_{0} - v_{r} \cdot \left( t-t_{0}\right)\right)
\label{eq:AbstandZuDeltaT}
\end{equation}

\subsection{Dauerstrichradar}
Das Dauerstrichradar, auch CW-Radar\footnote{\textbf{C}ontinuous  \textbf{W}ave Radar} sendet im Gegensatz zum Puls-Radar ein kontinuierliches Signal aus. Durch Laufzeitunterschiede des empfangenen und gesendeten Signals lässt sich so der Abstand zu einem Ziel bestimmen. In Abbildung\ref{fig:cwradar} ist das CW-Radar als Blockschaltbild dargestellt.
\begin{figure}[tbp]
  \centering
  \tikzset{%
	% Self defined bulding blocks. 
	% Nevertheless circutikz has implemented filters, couplers and other components since version 0.4, they are mostly implemented as bipoles.
	% The usage of bipoles: \draw (start) to[lowpass/amp/adc,....] (end).
	% The problem is, that if one wants to use arrows, the arrows in bipoles can not be sat manual (fixed in circuitikz source) AND THEY ARE NOT CONSISTENT
	% Also it's quite a mess, which component is a monopole, simple block, bipol, quad/triple etc
	% Following are a few examples on how to define your own blocks. 
	%
	% % % % % % % % % % % % % % % % % % % % % % % % % % % % % % % % % % % % % % % % % % % % % % % % % % % % % % % % % % % % % % % % % % % % %
	% % % % % % % % % % % % % % % % % % % % % % % % % % % % % % % % % % % % % % % % % % % % % % % % % % % % % % % % % % % % % % % % % % % % %
	% % % % % % % % % % % % % % % % % % % % % % % % % % % % % % % % % % % % % % % % % % % % % % % % % % % % % % % % % % % % % % % % % % % % %
	% % % % % % % % % % % % % % % % % % % % % % % % % % % % % % % % % % % % % % % % % % % % % % % % % % % % % % % % % % % % % % % % % % % % %
	%
	% Standard block definition, the width and height is adopted from the circutizk source code, so don't mind the strange values. Also the linewidth is set according to the circutrikz source code.
	block/.style    	= 	{draw, fill=white, thick, rectangle, minimum height = 0.98cm, minimum width = 0.98cm, node distance=2.5cm, line width=1.5pt},
	%
	% Standard circular block
	circleblock/.style	= 	{draw, fill=white, thick, circle, minimum width = 0.98cm,  line width=1.5pt, node distance=2.5cm},
	%
	% Label for circuitikz nodes, as they're reference is in the middle and not on the outer edge of the node....
	circuitikzlabel/.style	=	{label={[label, label distance=0.5cm]#1}},
	%
	%
	%
	% VCO/Oscillator 
	myVCO/.style			= 	{circleblock, path picture={%
		\draw[line width=0.75pt] 	($(path picture bounding box.west)+(0.09cm,0)$) sin ($(path picture bounding box.center)-(0.2cm,-0.2cm)$) cos  (path picture bounding box.center) sin ($(path picture bounding box.center)-(-0.2cm,0.2cm)$) cos ($(path picture bounding box.east)-(0.09cm,0)$);
		}
	},
	% Amplifier, as circuitikz does only provite amplifiers as 2-ports/bipoles
	myAMP/.style		= 	{block, node distance=2.5cm, path picture={%
		\draw[fill=white, line width=0.75pt] ($(path picture bounding box.center)+(0.7em,0)$) -- ($(path picture bounding box.center)-(0.7em,-0.7em)$) -- ($(path picture bounding box.center)-(0.7em,0.7em)$)  -- cycle;
		}
	},
	% Block	
	myBlock/.style    	= 	{draw, fill=white, thick, rectangle, minimum height = 0.98cm, minimum width = 0.98cm, node distance=2.5cm, line width=1.5pt},
	myBigBlock/.style    	= 	{draw, fill=white, thick, rectangle, minimum height = 0.98cm, minimum width = 2.94cm, node distance=2.5cm, line width=1.5pt},	
	% Same for ADC
	myADC/.style 	=	{block, path picture={%
		\draw[line width=0.75pt] 	(path picture bounding box.south west) -- (path picture bounding box.north east);
		\node[] at ($(path picture bounding box.center)+(-.5em,.5em)$) () {D};
		\node[] at ($(path picture bounding box.center)+(.5em,-.5em)$) () {A};
		} 
	},
	% Same for filters
	myLP/.style	=	{block, path picture={%
		%Sine-Waves
		\draw[line width=.75pt] 	($(path picture bounding box.west)+(0.3em,0)$) sin ($(path picture bounding box.center)-(0.50em,-0.3em)$) cos  (path picture bounding box.center) sin ($(path picture bounding box.center)-(-0.50em,0.3em)$) cos ($(path picture bounding box.east)-(0.3em,0)$);
		\draw[line width=0.75pt] 	($(path picture bounding box.west)+(0.3em,-0.65em)$) sin ($(path picture bounding box.center)-(0.50em,0.35em)$) cos  ( $(path picture bounding box.center)-(0,0.65em)$) sin ($(path picture bounding box.center)-(-0.50em,0.95em)$) cos ($(path picture bounding box.east)-(0.3em,0.65em)$);
		\draw[line width=0.75pt] 	($(path picture bounding box.west)+(0.3em,0.65em)$) sin ($(path picture bounding box.center)-(0.50em,-0.95em)$) cos  ( $(path picture bounding box.center)+(0,0.65em)$) sin ($(path picture bounding box.center)-(-0.50em,-0.35em)$) cos ($(path picture bounding box.east)-(0.3em,-0.65em)$);
		% Cancelation
		\draw[line width=0.75pt] 	($(path picture bounding box.center)-(0.2em,0.2em)$) -- (path picture bounding box.center) -- ($(path picture bounding box.center)+(0.2em,0.2em)$) ;
		\draw[line width=0.75pt] 	($(path picture bounding box.center)-(0.2em,-0.45em)$) -- ($(path picture bounding box.center)+(0,0.65em)$) -- ($(path picture bounding box.center)+(0.2em,0.85em)$) ;
		}
	},
}
\begin{tikzpicture}[line width=0.7pt,>=latex,node distance=2.5cm]
	% First: All building blocks are placed relative to the first component
	\draw (0,0)
		node[myVCO, label={above:Signalgenerator}] (oszi) {}
		% As the coupler ports are not in the middle, based on the size (again extraceted from circutikz source code), an yshift is perfomed to have the input on the same height as the output of the oszillator. The xshift is used to place the VCO and ADC on the same y-value after all.
		% Undo the yshift as the output of the coupler is on the same height as the input of the amplifier
		node[myBlock, right of=oszi,  node distance=5cm, label={above:Sender}] (pa) {$P_{S}$}
		% Circulator is rotated that the ports are on the correct position, normally ports are arranged as follows:
		% 
		%				  o------------o
		%						|
		%						o
		node[circulator, below right of=pa, xshift = 3cm, rotate=90,] (circ) {}
		node[antenna, right of = circ] (antenna) {}
		node[myBlock, below left of=circ,  label={below:Empfänger}] (lna) {$P_{E}$}
		% Used redefinition of label, otherwise the label would be overlapping with the mixer shape
		node[mixer, below of=pa,yshift = -1cm  , circuitikzlabel={below:Mischer}] (mixer) {}
		node[myLP, left of=mixer, label={below:Tiefpass}] (lowpass) {}
		node[myADC, left of = lowpass,node distance=2.5cm, label ={below:Signalverarbeitung}](Sig){}		
		;
		
	% Connect everything together
	\draw[->] (oszi) -- node[above]{$u_{S}(t)$}(pa.west);

	%\draw[] (coupler.1) -| ++(-0.5cm,-0.1cm) node[match, rotate=-90, yscale=-1] {};
	\draw[->] (pa.east) -| node[above left]{$u_{S}(t)$}(circ.2);
	\draw[-] (circ.3) -- (antenna);
	\draw[->] (circ.1) |- node[above left]{$u_{E}(t)$}(lna.east);
	\draw[->] (lna.west) -- node[above]{$u_{E}(t)$}(mixer.east);
	\draw[->] (mixer.west)  -- node[above]{$u_{M}(t)$} (lowpass.east);
	\draw[->] (lowpass.west) -- node[above]{$u_{TP}(t)$}(Sig.east);
	\draw[->] (pa.south) -- node[right]{$u_{S}(t)$}(mixer.north);

\end{tikzpicture}

  \caption{CW-Radar Blockdiagramm}
  \label{fig:cwradar}
\end{figure}
Das Sendesignal $u_{S}(t)$ wird versendet. Das empfangene Signal $u_{E}(t)$ gelangt in den Mischer. Im Mischer wird das empfangene Signal mit dem ursprünglichen Signal gemischt. Daraus ergibt sich das Mischsignal $u_{M}(t)$. Das Mischsignal wird in dem nachfolgenden Tiefpass gefiltert. Das daraus entstehende Tiefpasssignal $u_{TP}(t)$ wird weiterverarbeitet und digitalisiert. \\
Ist das Sendesignal $u_{S}(t)$ mit
\begin{align}
u_{S}(t) =  u_{S} \cos\left( \omega_{S}t + \varphi_{S}\right)
\end{align} 
cosinus-förmig mit Phase $\varphi_{S}$ und Winkelgeschwindigkeit $\omega_{S}$ mit 
\begin{align}
\omega_{S} = 2 \pi f_{s}
\end{align}
und das Empfangssignal  $u_{E}(t)$ 
\begin{align}
u_{E}(t) =  u_{E} \cos\left( \omega_{S}\left(t-\Delta t\right) + \varphi_{S}+\varphi_{E}\right) 
\end{align}
ergibt sich im Mischer ein Mischsignal $u_{M}(t)$ mit
\begin{align}
u_{M}(t) = u_{S}(t)\cdot u_{E}(t) = u_{S} u_{E} \cos\left( \omega_{S}t + \varphi_{S}\right)\cdot \cos\left( \omega_{S}\left(t-\Delta t\right) + \varphi_{S}+\varphi_{E}\right).
\end{align}
Wird die trigonometrische Beziehung 
\begin{align}
\cos \alpha \cdot \cos \beta = \frac{1}{2} \left( \cos \left( \alpha + \beta \right) + \cos \left( \alpha - \beta \right) \right)
\end{align}
betrachtet, ergibt sich für das Mischsignal zwei Frequenzen. Zum Einen die Summe der Frequenzen $2\omega_{s}$ zum Anderen die Subtraktion $-\omega_{s}$. Der Tiefpass wird darauf ausgelegt, dass 
\begin{align}
f_{g,TP} = 2f_{s}
\end{align}
gilt.
Es ist nach \cite[S.43,S.44]{HuderRadar} das Tiefpass-Ausgangssignal $u_{TP}(t)$
\begin{align}
u_{TP}(t) = u_{TP} \cos \left( -\omega_{s} \Delta t + \varphi_{Z}\right)
\end{align} und mit Gleichung \ref{eq:AbstandZuDeltaT} ist
\begin{align}
u_{TP}(t) = u _{TP} \cdot \cos \left( \frac{2 \omega_{S} }{c_{0}}\cdot \left( -R_{0} + v_{r}\cdot t - 2 v_{r}\cdot t_{0}\right) +\varphi_{E}\right). 
\end{align}
Ist die Radialgeschwindigkeit $v_{r}$ gleich Null, also keine Bewegung der Sender/Empfängersysteme oder des Ziels, ergibt sich ein konstantes Ausgangssignal am Tiefpass. Dieses konstante Signal lässt sich in die Entfernung zum Ziel umrechnen. 
\subsection{Dauerstrichradar mit Frequenzmodulation}
Das Dauerstrichradar mit Frequenzmodulation, auch FMCW\footnote{\textbf{F}requency \textbf{M}odulated \textbf{C}ontinous \textbf{W}ave}-Radar genannt, beschreibt ein Dauerstrichradar, in welchem das Sendesignal frequenzmoduliert und nicht amplitudenmoduliert wird. Dabei ist der gerätetechnische Aufbau in Abbildung \ref{fig:FMCWRadar} dargestellt. 
\begin{figure}[tbp]
  \centering
  \input{gfx/FMCW_Block_Bjoern.tikz}
  \caption{Blockschaltbild FMCW-Radar}
  \label{fig:FMCWRadar}
\end{figure}
Das aus dem Frequenzmodulator frequenzmodulierte Signal wird über den Sender abgestrahlt. Dabei ergibt sich für das Sendesignal
\begin{align}
u_{s}(t) = u_{s} \cdot \cos\left( 2\pi f_{s}(t) + \varphi_{s} \right) 
\end{align}     
mit Amplitude $u_{s}$ und Phase $\varphi_{s}$. Durch den Frequenzmodulator ergibt sich zusätzlich für die Sendefrequenz $f_{s}(t)$
\begin{align}
f_{s}(t) = f_{0} + \frac{\Delta f}{T_{m}}\cdot t.
\end{align}
\begin{figure}[tbp]
  \centering
  \tikzset{%
	% Self defined bulding blocks. 
	% Nevertheless circutikz has implemented filters, couplers and other components since version 0.4, they are mostly implemented as bipoles.
	% The usage of bipoles: \draw (start) to[lowpass/amp/adc,....] (end).
	% The problem is, that if one wants to use arrows, the arrows in bipoles can not be sat manual (fixed in circuitikz source) AND THEY ARE NOT CONSISTENT
	% Also it's quite a mess, which component is a monopole, simple block, bipol, quad/triple etc
	% Following are a few examples on how to define your own blocks. 
	%
	% % % % % % % % % % % % % % % % % % % % % % % % % % % % % % % % % % % % % % % % % % % % % % % % % % % % % % % % % % % % % % % % % % % % %
	% % % % % % % % % % % % % % % % % % % % % % % % % % % % % % % % % % % % % % % % % % % % % % % % % % % % % % % % % % % % % % % % % % % % %
	% % % % % % % % % % % % % % % % % % % % % % % % % % % % % % % % % % % % % % % % % % % % % % % % % % % % % % % % % % % % % % % % % % % % %
	% % % % % % % % % % % % % % % % % % % % % % % % % % % % % % % % % % % % % % % % % % % % % % % % % % % % % % % % % % % % % % % % % % % % %
	%
	% Standard block definition, the width and height is adopted from the circutizk source code, so don't mind the strange values. Also the linewidth is set according to the circutrikz source code.
	block/.style    	= 	{draw, fill=white, thick, rectangle, minimum height = 0.98cm, minimum width = 0.98cm, node distance=2.5cm, line width=1.5pt},
	%
	% Standard circular block
	circleblock/.style	= 	{draw, fill=white, thick, circle, minimum width = 0.98cm,  line width=1.5pt, node distance=2.5cm},
	%
	% Label for circuitikz nodes, as they're reference is in the middle and not on the outer edge of the node....
	circuitikzlabel/.style	=	{label={[label, label distance=0.5cm]#1}},
	%
	%
	%
	% VCO/Oscillator 
	myVCO/.style			= 	{circleblock, path picture={%
		\draw[line width=0.75pt] 	($(path picture bounding box.west)+(0.09cm,0)$) sin ($(path picture bounding box.center)-(0.2cm,-0.2cm)$) cos  (path picture bounding box.center) sin ($(path picture bounding box.center)-(-0.2cm,0.2cm)$) cos ($(path picture bounding box.east)-(0.09cm,0)$);
		}
	},
	% Amplifier, as circuitikz does only provite amplifiers as 2-ports/bipoles
	myAMP/.style		= 	{block, node distance=2.5cm, path picture={%
		\draw[fill=white, line width=0.75pt] ($(path picture bounding box.center)+(0.7em,0)$) -- ($(path picture bounding box.center)-(0.7em,-0.7em)$) -- ($(path picture bounding box.center)-(0.7em,0.7em)$)  -- cycle;
		}
	},%%
	% Block	
	myBlock/.style    	= 	{draw, fill=white, thick, rectangle, minimum height = 0.98cm, minimum width = 0.98cm, node distance=2.5cm, line width=1.5pt},
	myBigBlock/.style    	= 	{draw, fill=white, thick, rectangle, minimum height = 0.98cm, minimum width = 2.94cm, node distance=2.5cm, line width=1.5pt},	
	% Same for ADC
	myADC/.style 	=	{block, path picture={%
		\draw[line width=0.75pt] 	(path picture bounding box.south west) -- (path picture bounding box.north east);
		\node[] at ($(path picture bounding box.center)+(-.5em,.5em)$) () {D};
		\node[] at ($(path picture bounding box.center)+(.5em,-.5em)$) () {A};
		} 
	},
	% Same for filters
	myLP/.style	=	{block, path picture={%
		%Sine-Waves
		\draw[line width=.75pt] 	($(path picture bounding box.west)+(0.3em,0)$) sin ($(path picture bounding box.center)-(0.50em,-0.3em)$) cos  (path picture bounding box.center) sin ($(path picture bounding box.center)-(-0.50em,0.3em)$) cos ($(path picture bounding box.east)-(0.3em,0)$);
		\draw[line width=0.75pt] 	($(path picture bounding box.west)+(0.3em,-0.65em)$) sin ($(path picture bounding box.center)-(0.50em,0.35em)$) cos  ( $(path picture bounding box.center)-(0,0.65em)$) sin ($(path picture bounding box.center)-(-0.50em,0.95em)$) cos ($(path picture bounding box.east)-(0.3em,0.65em)$);
		\draw[line width=0.75pt] 	($(path picture bounding box.west)+(0.3em,0.65em)$) sin ($(path picture bounding box.center)-(0.50em,-0.95em)$) cos  ( $(path picture bounding box.center)+(0,0.65em)$) sin ($(path picture bounding box.center)-(-0.50em,-0.35em)$) cos ($(path picture bounding box.east)-(0.3em,-0.65em)$);
		% Cancelation
		\draw[line width=0.75pt] 	($(path picture bounding box.center)-(0.2em,0.2em)$) -- (path picture bounding box.center) -- ($(path picture bounding box.center)+(0.2em,0.2em)$) ;
		\draw[line width=0.75pt] 	($(path picture bounding box.center)-(0.2em,-0.45em)$) -- ($(path picture bounding box.center)+(0,0.65em)$) -- ($(path picture bounding box.center)+(0.2em,0.85em)$) ;
		}
	},
}
\begin{tikzpicture}[line width=0.7pt,>=latex,node distance=2.5cm,scale = 1.2]

	%Achsen oben
	\draw[->] (-0.5,0) -- (8,0);
	\draw[->] (0,-0.5) -- (0,4); 
	\draw(8,0) node[label={below: $t$}](){};
	\draw(0,4) node[label={left: $f(t)$}](){};
	
	%Achsen unten
	\draw[->] (-0.5,-3)--(8,-3);
	\draw(0,-1) node[label={left: $f_{\text{ZF}}(t)$}](){};
	\draw[->] (0,-3.5)--(0,-1);
	\draw(8,-3) node[label={below: $t$}](){};
	
	%Dashed oben
	\draw[dashed](-0.25,3.5) --(8,3.5);
	\draw(-0.25,3.5) node[label={left: $f_{0}$}](){};
	\draw[dashed](-0.25,0.5) --(8,0.5);
	\draw(-0.25,0.5) node[label={left: $f_{0} + \Delta f$}](){};
	
	%Label Bereich Delta t
	\draw[-] (0.8,0)--(0.8,-0.5);
	\draw[<->](0,-0.3) -- node[below]{$\Delta t$}(0.8,-0.3);
	
	%Label Rechts Delta f 
	\draw[<->](7.8,0.5)--node[right]{$\Delta f$}(7.8,3.5);
	
	%fzf
	\draw[<->](2,1.7)--(2,2.5);
	\draw (2,2.5) node[label ={ left :$f_{\text{ZF}}(t)$}](){};

	%Schwarz unten
	\draw[-] (0,-1.5) -- (0.8,-1.5);
	\draw[-] (0.8,-1.5) -- (0.8,-2.5);
	\draw[-] (0.8,-2.5) -- (3,-2.5);
	\draw[-] (3,-2.5) -- (3,-1.5);
	\draw[-] (3,-1.5) -- (3.8,-1.5);
	\draw[-] (3.8,-1.5) --(3.8,-2.5);
	\draw[-] (3.8,-2.5) --(6,-2.5);
	\draw[-] (6,-2.5) -- (6,-1.5);
	\draw[-] (6,-1.5) --(6.8,-1.5);
	\draw[-] (6.8,-1.5) --(6.8,-2.5);
	\draw[-] (6.8,-2.5) --(7.5,-2.5);
		
	%Dashed Rot
	\draw[dashed,red](0,2.7) -- (0.8,3.5);
	\draw[dashed,red](0.8,3.5) -- (0.8,0.5);
	\draw[dashed,red](0.8,0.5) -- (3.8,3.5);
	\draw[dashed,red](3.8,3.5) -- (3.8,0.5);
	\draw[dashed,red](3.8,0.5) -- (6.8,3.5);
	\draw[dashed,red](6.8,3.5) -- (6.8,0.5);
	\draw[dashed,red](6.8,0.5) -- (7.5,1.2);
	
	% Blau
	\draw[-,blue](0,0.5)--(3,3.5);
	\draw[-,blue](3,3.5)--(3,0.5);
	\draw[-,blue](3,0.5)--(6,3.5);
	\draw[-,blue](6,3.5)--(6,0.5);
	\draw[-,blue](6,0.5)--(7.5,2);
	
	%Achsen Tm
	\draw[-](3,0.1)--node[below]{$T_{m}$}(3,-0.1);
	\draw[-](3,-2.9)--node[below]{$T_{m}$}(3,-3.1);	
	
	%Legende
	\draw[blue,-] (1,4.5)--node[right]{Sendesignal $f_{\text{S}}(t)$}(1.3,4.5);
	\draw[red,dashed] (4,4.5)--node[right]{Empfangssignal $f_{\text{E}}(t)$}(4.3,4.5);
\end{tikzpicture}

  \caption{Frequenzverlauf Sende-,  Empfangs-und Zwischenfrequenzsignal\cite[S.6]{Huegler}\cite[S.78]{HuderRadar}}
  \label{fig:FMCW_Frequenz}
\end{figure}
Dabei sind in Abbildung\ref{fig:FMCW_Frequenz} der Frequenzhub $\Delta f$, die Modulationsperiode $T_{m}$ sowie der Frequenzoffset $f_{0}$ dargestellt. Der Frequenzverlauf ist hier als sägezahnförmig angenommen. Andere Formen, beispielsweise sinuförmig oder linear, sind ebenfalls möglich. Zusätzlich wird auf die erste Modulationsperiode beschränkt. Für weitere Perioden gilt das Vorgehen analog.\\
Das Empfangssignal $s_{e}(t)$ wird um den Laufzeitunterschied $\Delta t$ mit
\begin{align}
\Delta t = \frac{2R}{c_{0}}
\end{align}
verschoben. Dabei ist R der Radar-Ziel Abstand. Dadurch ergibt sich für die Empfangsfrequenz $f_{e}(t)$ 
\begin{align}
f_{e}(t) = f_{0} + \frac{\Delta f}{T_{m}}\cdot \left( t-\Delta t \right).
\end{align}
Ist zudem der in Abbildung \ref{fig:FMCWRadar} dargestellte Tiefpass darauf ausgelegt, die Summenfrequenz $f_{s} + f_{e}$ nicht passieren zu lassen, ergibt sich analog zu \cite[S.43,S.44]{HuderRadar} für die Zwischenfrequenz
\begin{align}
f_{ZF}(t) = \vert f_{s}(t)- f{e}(t) \vert = \frac{\Delta f}{T_{m}}\cdot \Delta T.
\end{align}
Damit ist die Zwischenfrequenzspannung $u_{ZF}(t)$ 
\begin{align}
u_{ZF}(t) = u_{ZF} \cdot  \cos\left( 2\pi \frac{\Delta f \Delta t}{T_{m}} \cdot t + \varphi_{s} \right) .
\end{align}
Daraus ergibt sich nach\cite[S.80]{HuderRadar} die Zielentfernung $R$ zu
\begin{equation}
R = \frac{c_{0}T_{m}}{2\Delta f}\cdot f_{ZF}.
\end{equation}
Um $R$ zu bestimmen muss die Zwischenfrequenz gemessen werden. Dafür wird das Zwischenfrequenzsignal verarbeitet und anschließend digitalisiert. Das Spektrum des Digitalsignals kann mit einer FFT\footnote{\textbf{F}ast \textbf{F}ourier \textbf{T}ransformation} berechnet werden. Daraus ergibt sich ein diskretes Spektrum, in dem die Zwischenfrequenz erkennbar ist. Falls mehrere Ziele detektiert werden, lassen sich ebenso mehrere Zwischenfrequenzen erkennen.
\subsection{Chirp-Sequence-Radar}
Ein Nachteil des Dauerstrichradar mit Frequenzmodulation ist, dass es nicht ohne weiteres möglich ist, die Radialgeschwindigkeit zu bestimmen. Dafür kann das Dauerstrichradar mit Frequenzmodulation erweitert werden zu einem Chirp-Sequence-Radar. Dabei werden viele, möglichst kurze Rampen als Frequenzmodulation vorgenommen - die einzelnen Rampen sind demnach möglichst Steil. Das Chirp-Sequence-Radar wird zudem nicht über eine einzelne Rampe ausgewertet, sondern über eine feste Menge an Rampen. Im Spektrum entstehen nun eindeutig voneinander trennbare spektrale Anteile die entweder proportional zum Abstand oder zur Geschwindigkeit sind. Bei ausreichender Steilheit finden sich die Anteile die proportional zum Abstand sind in den oberen Bereichen des Spektrums, entsprechend die Anteile der Geschwindigkeit in unteren Teilen des Spektrums. 
\chapter{Signalkonditionierung}
Signalkonditionierung beschreibt die Bearbeitung eines analogen Signals vor der Weiterverarbeitung in der nächsten, oft letzten, Systemstufe. Ziel der in dieser Arbeit durchgeführten Signalkonditionierung ist die Aufbereitung der vom Radarchip empfangenen Zwischenfrequenzsignale für die Weiterleitung in einen Analog-Digital-Konverter. Dafür sollen benötigte Systemkomponenten vorgestellt und mit den wichtigsten Systemparametern beschrieben werden. 
\section{Rauschen}
Rauschen beschreibt eine stochastische Störgröße die sich in der Änderung von Signalen zeigt. Dabei begrenzt Rauschen die Möglichkeit von Systemen, Signale eindeutig zu detektieren und muss demnach stets mit betrachtet werden wenn ein System oder ein Kommunikationsweg ausgelegt wird. Dabei ist insbesondere das Verhältnis von Signalleistung zu Rauschleistung, das SNR, entscheidend. Dies ist über
\begin{align}
\text{SNR} = \frac{S}{N} = \frac{P_{Signal}}{P_{Rauschen}}
\end{align}
definiert. Dabei wird das SNR meist in dB angegeben. Ursachen von Rauschen sind unterschiedlich. Folgend sind einige wichtige Rauschprozesse aufgezählt.
\begin{description}
\item Thermisches Rauschen\\
Thermisches Rauschen wird durch die Behinderung freier Ladungsträger durch Gitterschwingungen in festen Leitern verursacht. Dabei verändert sich die Geschwindigkeit der Ladungsträger und erzeugt somit das Thermische Rauschen. Auf thermisches Rauschen wird weiterhin noch genauer eingegangen.
\item Schrotrauschen \\
Schrotrauschen wird durch diskrete Ladungsträger verursacht. Dabei ist bei konstantem Strom die Anzahl der Ladungsträger nicht konstant. Dies wird durch verschiedene energetische Barrieren oder Rekombinationsprozesse in Halbleiterübergängen verursacht. Da jeder Ladungsträger, im Fall von elektrischen Strom Elektronen, zum Rauschen beiträgt, lässt sich das gemittelte Quadrat des Rauschstroms $\bar{i}_{n}^2$ über
\begin{align}
\bar{i}_{n}^2 = 2eI\Delta f
\end{align}
mit I dem Leiterstrom, e der Elementarladung und $\Delta f $ der Messbandbreite, angeben\cite{Rauscherino}. Es ist auffällig, dass Schrotrauschen nicht direkt von der Frequenz, sondern von dem Frequenzband abhängt.
\item 1/f Rauschen \\
1/f Rauschen entsteht aus den Eigenschaften von Grenzflächenzuständen\cite{HFKompo}.  Dabei nimmt die Rauschleistung mit 1/f ab. Ab etwa \SI{10}{k\hertz} ist dieser Rauschtyp vernachlässigbar klein.
\end{description}

\subsection{Thermisches Rauschen}
Es wurde bereits der Hintergrund von thermischen Rauschen genannt. Dabei kann bis zu hohen Frequenzen angenommen werden, dass thermisches Rauschen frequenzunabhängig ist. So ist die thermische Rauschleistung eines Widerstands $P_{n}$ mit
\begin{align}
P_{n} = \frac{\bar{u}_{n}^2}{4R} = kT\Delta f
\end{align}
gegeben. Stellt man nach $\bar{u}_{n}^2$ um, ergibt sich für den Mittelwert der Rausch-Spannung
\begin{align}
\bar{u}_{n}^2 = 4kTR\Delta f.
\end{align}
Analog lässt sich somit die Leistungsdichte der Rauschspannung nach \cite{Rauscherino} zu
\begin{align}
S_{n} = \frac{\bar{u}_{n}^2}{\Delta f} = 4kTR
\end{align}
bestimmen. Dabei ist $S_{n}$ über die Frequenz konstant. \\
Eine Möglichkeit um einen rauschenden Widerstand zu beschreiben ist, ihn in einen idealen, nicht rauschenden Widerstand sowie eine rauschende Quelle aufzuteilen. Dabei kann entweder eine Rausch-Stromquelle oder eine Rausch-Spannungsquelle gewählt werden. Für eine Rausch-Stromquelle wird der Kurzschlussstorm
\begin{align}
\bar{i}_{n,KS}^2 = \frac{4kT\Delta f}{R}
\end{align}
und für eine Rausch-Spannungsquelle die Leerlaufspannung
\begin{align}
\bar{u}_{n,LL}^2 = 4kTR\Delta f 
\end{align}
bestimmt\cite{HFKompo}.
\subsection{Phasenrauschen}
Phasenrauschen tritt stets dann auf, wenn ein periodisches (Zeit-)Signal erzeugt werden soll. Phasenrauschen bezeichnet die dabei auftretenden Effekte im Frequenzbereich und Jitter entsprechend die Korrelation im Zeitbereich. Ein periodisches Zeitsignal korreliert mit einem diskreten Spektrum, einem Dirac-Impuls(Bild \ref{Oszi} links). In der Realität sind diskrete Signale unmöglich zu erreichen. Ein realer Dirac-Impuls ist dabei in Bild \ref{Oszi} rechts dargestellt. Die benachbarten Frequenzanteile neben $f_{0}$ werden als Phasenrauschen bezeichnet. Um Phasenrauschen zu senken, werden Phasenregelschleifen verwendet.\\
\begin{figure}[tbp]
  \centering
  \begin{tikzpicture}[scale = 1.3]
	%x-Achse 1
	\draw[->] (-0.5,0) -- (4,0) node[anchor=north] {$f$};
	%label f0
	\draw (2,0) node[anchor=north]{$f_{0}$};
	%y-Achse 1
	\draw[->] (0,-0.5) -- (0,4) node[anchor=east] {Leistungsdichte};
	
	
	%x-Achse 2
	\draw[->] (5.5,0) -- (10,0) node[anchor=north] {$f$};
	%label f0
	\draw (8,0) node[anchor=north]{$f_{0}$};
	%y-Achse 2
	\draw[->] (6,-0.5) -- (6,4) node[anchor=east] {Leistungsdichte};
	
	%Dirac Links
	\draw[-] (2,0) -- (2,3);
	
	%Dirac Dottet Rechts
	\draw[dotted] (8,0) -- (8,3);
	
	%FrequenzSeitenBänder
	\draw[-] (8,3) parabola[bend at end](10,0);
	\draw[-] (8,3) parabola[bend at end](6,0);
	
\end{tikzpicture}

  \caption{Gegenüberstellung Spektrum ideal(links) und real(rechts) Oszillators nach \cite{HFKompo}}
  \label{Oszi}
\end{figure}

Um die Oszillator-Stabilität zu bestimmen, muss zusätzlich zwischen Kurzzeitstabilität und Langzeitstabilität unterschieden werden. \\
Die Langzeitstabilität wird insbesondere durch die Temperatur oder Alterungsprozssse beeinflusst. Dabei ändert sich die Mittenfrequenz $f_{0} $ des Oszillators. Die Kurzzeitstabilität bewirkt die Ausbreitung von Rauschseitenbändern. Dies wird durch momentane stochastische, periodische oder deterministische Schwankungen zurückgeführt\cite{HFKompo}.
\section{Verstärker}
Ein Verstärker bezeichnet eine Systemkomponente, welche ein  analoges Eingangsignal, in Form von Strom, Spannung oder Leistung, in ein analoges Ausgangssignal nicht zwangsläufig gleicher Art überführt und dieses aktiv verstärkt. Dabei ist die wichtigste Kenngröße die Verstärkung oder auch der Gain. Dabei wird der Gain stets über 
\begin{center}
Ausgangssignal $ = $ Eingangssignal $ \cdot $ Gain
\end{center}
beschrieben. Häufig wird Gain $G$ in dB angegeben, dabei gilt
\begin{equation}
\frac{G}{\text{dB}} = \text{B} \cdot \log \left( \frac{G}{\text{lin}} \right).
\end{equation}
Dabei ist 
\begin{equation}
\text{B} = \begin{cases} 20 \text{ , falls Signal Spannung oder Storm} \\ 10 \text{ , falls Signal Leistung} \end{cases}.
\end{equation}
Dabei unterscheiden sich verschiedene Verstärkertypen. Einige sind folgend aufgezählt.
\begin{description}
\item Stormverstärker \\
	  Eingangsstrom in Ausgangsstorm
\item Spannungsverstärker    		\\
	  Eingangsspannung in Ausgangsspannung
\item Transkonduktanzverstärker   \\
	  (Eingangsspannung  in Ausgangsstrom
\item Transimpedanzverstärker  	\\
Eingangsstrom  in Ausgangsspannung
\end{description}
Zusätzlich lassen sich Verstärker bezüglich ihrer Funktion oder Verwendung  charakterisieren. Es wird sich weiterhin auf Operationsverstärker beschränkt. In Abbildung \ref{fig:OpAmp} ist ein solcher dargestellt. Operationsverstärker haben einen differentiellen Eingang und verstärken Spannungen, sind also ein Spannungsverstärker. Dabei soll die Differenz der Eingänge $U_{in,p}$ und $U_{in,n}$ die Ausgangsspannung $U_{aus}$ bestimmen. Also gilt
\begin{equation}
U_{aus} = G \cdot \left( U_{in,p} - U_{in,n}\right).
\end{equation} Dabei hat der Operationsverstärker meist nur einen Ausgang. Differentielle Ausgänge werden im Anschluss durch invertieren eines Eingangs erzeugt. Invertiert wird durch einen Operationsverstärker mit $G = -1$ sowie dem negativem Input auf dem Bezugspotenzial.\\ 
Ein Operationsverstärker benötigt eine eigene Versorgungsspannung $V_{S}$. Diese gibt die maximale Ausgangsspannung an. Die Eingangsimpedanz des Operationsverstärkers sollte hoch sein. In idealen Operationsverstärkern wird sie als $\text{Z}_{in} = \inf$ angenommen. Dagegen ist die Ausgangsimpedanz möglichst klein, in idealen Operationsverstärkern wird diese als $\text{Z}_{aus} = 0$ angenommen. Also sind beide einfließenden Ströme $I_n = I_p = 0$.
\begin{figure}[tbp]
  \centering
  \begin{tikzpicture}[scale=2]
        \draw
        (0, 0) node[op amp,scale = 2] (opamp) {}
		(opamp.+) node[left] {$U_{in,n}$}
		(opamp.-) node[left] {$U_{in,p}$}
		(opamp.out) node[right] {$U_{aus}$}
		(opamp.up) --++(0,0.5) node[vcc]{$V_{S}$}
		(opamp.down) --++(0,-0.5) node[vee]{$-V_{S}$
}		     
        ;
\end{tikzpicture}

  \caption{Symbol eines Operationsverstärker}
  \label{fig:OpAmp}
\end{figure}
\subsection{Systemparameter von Operationsverstärkern}
Reale Operationsverstärker haben einige nicht ideale Effekte. Um diese zu beschreiben, werden Operationsverstärker mit einigen Parametern beschrieben
\begin{description}
\item Endlicher Gain \\
Dabei gilt es zwischen Open-Loop-Gain und Closed-Loop-Gain zu unterscheiden. Der Open-Loop-Gain bezeichnet die mögliche Verstärkung der Operationsverstärkers, wenn keine Rückführung auf die Eingänge vorhanden ist. Der Open-Loop-Gain eines idealen Operationsverstärkers ist $\inf$. Auch in realen Operationsverstärkern kann dieser Wert sehr groß werden. Der Closed-Loop-Gain wird durch die Beschaltung der Rückführung bestimmt sowie die Verhältnisse der dortigen Komponenten, meist Widerstände, definiert.
\item Eingangsimpedanz \\
Die differentielle Eingangsimpedanz gibt die Impedanz zwischen positiven und  negativen Eingang des Operationsverstärkers an. Die Gleichtakt Eingangsimpedanz, auch Common-Mode Impedanz, beschreibt die jeweilige Impedanz bezogen auf das Massepotenzial.
\item Gleichtaktunterdrückung \\
Gleichtaktunterdrückung CMR\footnote{ \textbf{C}ommon \textbf{M}ode \textbf{R}ejection} beschreibt die Änderung von $U_{aus}$ wenn sich beide Eingänge um den gleichen Wert ändern. Ein wesentliche Maß für die Qualität eines Operationsverstärkers ist die Gleichtaktunterdrückung. Dabei soll diese möglichst groß, idealerweise unendlich groß sein. Entsprechend ist die Gleichtaktverstärkung idealerweise null. 
\item Slew Rate und Bandbreite \\
Da ein Operationsverstärker stets seinen Eingängen folgt, ist dies nur bis zu einer maximalen Geschwindigkeit möglich. Diese Geschwindigkeit ist die Slew Rate. Ist die Änderung schneller als der Operationsverstärker, wird nicht mehr der volle Gain erreicht. Sobald der Gain um \SI{3}{dB} kleiner ist als vorgesehen, ist die Bandbreite überschritten.
\item Ausgangsrauschen\\
Da Operationsverstärker aktive Bauelemente sind, erzeugen diese stets auch ein eigenes Rauschen welches dem bisherigen Rauschpegel hinzugefügt wird. Ein wesentliches Qualitätsmerkmal eines Operationsverstärkers ist dabei das Ausgangsrauschen. Ein idealer Operationsverstärker hat dabei kein Ausgangsrauschen. In der Realität hängt es allerdings von Lastwiderstand sowie Gain ab. 
\item Interception Point 3 \\
Interception Points modellieren die in der Realität vorhandenen Nichtlinearitäten von Operationsverstärkern mit Hilfe von Polynomen niedriger Ordnung, beispielsweise einer Taylorentwicklung. Der Interception Point 3 bezeichnet demnach den Einfluss der nichtlinearen Terme dritter Ordnung. Dabei ist der Interception Point 3 keine physikalische Größe. Effektiv beschreibt der Interception Point 3 die dynamische Reichweite bzw. die Grenze davon. 
\end{description}
\subsection{Rauschen in Operationsverstärkern}
Um Verstärkerrauschen zu betrachten, wird ein Ersatzschaltbild für einen Verstärker samt Rückkopplung dargestellt\ref{fig:Verstaerkerrauschen_Block}. Dabei ist $e^{\text{j}\varphi}$ das ungewollte Phasenrauschen. Die Übertragungsfunktion der Rückkopplung $\beta(\omega)$ ergibt sich aus der Schwingfrequenz des Oszillators.
\begin{figure}[tbp]
  \centering
  \input{gfx/Verstaerkerrauschen_Block.tikz}
  \caption{Blockschaltbild Zweitor Oszillator mit Verstärkung G \cite{HFKompo}}
  \label{fig:Verstaerkerrauschen_Block}
\end{figure}
Nach dem Barkhausen-Kriterium\cite[S.134,135]{HFKompo} liefert, dass der Betrag der Schleifenverstärkung gleich eins sein muss. Demnach gilt
\begin{align}
\vert \beta ( \omega )\cdot G\vert = 1 .
\end{align}
Zusätzlich muss die Phasenverschiebung zwischen Eingang und Ausgang 
\begin{align}
\text{arg}\left( \beta ( \omega )\right) + \varphi = 2\pi n
\end{align}
mit n aus den natürlichen Zahlen inklusive der Null. \\
Das Verstärkerrauschen ergibt sich aus der Summe der einzelnen Rauschanteile. Aus thermischen Rauschen, welches Frequenzunabhängig ist, ergibt sich die spektrale Leistungsdichte der Phasenschwankung \cite{HFKompo} 
\begin{align}
S_{\varphi , Thermisch} = 2 G \frac{kFT_{0}}{P_{0}} 
\end{align} 
für kleine Phasenschwankungen. Zudem ergibt sich aus dem 1/f Rauschen ein mit 1/f abfallender Anteil für niedrige Frequenzen. Daraus ergibt sich die spektrale Rauschdichte eines Verstärkers zu
\begin{align}
S_{\varphi}(\delta f) = \frac{\alpha_{1/f}}{(\delta f)^{-1}} + \frac{S_{\varphi , Thermisch}}{(\delta f)^{0}}.
\end{align}
Dabei ist $\alpha_{1/f}$ eine Konstante des 1/f Rauschens. In Bild \ref{fig:Verstaerkerrauschen} ist der Verlauf der Rauschdichte dargestellt.
\begin{figure}[tbp]
  \centering
  \begin{tikzpicture}[scale = 1.3]
	%x-Achse 1
	\draw[->] (-0.5,0) -- (6,0) node[anchor=north] {$\delta f (log)$};
	%label f0
	\draw (2,0) node[anchor=north]{$f_{0}$};
	%y-Achse 1
	\draw[->] (0,-0.5) -- (0,4) node[anchor=east] {$S_{\varphi} (\delta f)$ (log)};
	
	\draw (0.5,3) -- node[above right]{1/f Rauschen}(3.5,1);
	\draw (3.5,1) -- node[above]{Thermisches Rauschen}(6,1);

	
\end{tikzpicture}

  \caption{Verstärkerrauschen nach \cite{HFKompo}}
  \label{fig:Verstaerkerrauschen}
\end{figure}
\section{Filter}
Ein Filter ist eine Systemkomponente die die Amplitude und Phase eines Signals in Bezug in Abhängigkeit der Frequenz verändert. Dabei können Filter aktive oder passive Bauelemente sein, und auch abweichende Eigenschaften von den bereits genannten besitzen. Ein Tiefpassfilter lässt tiefe Frequenzen passieren und dämpft hohe. Hochpassfilter dämpfen entsprechend tiefe Frequenzen und lässt hohe passieren.
\subsection{Tiefpassfilter}
Ein Tiefpassfilter besteht aus einem Durchlass- und einem Sperrbereich. Das Filterverhalten in den jeweiligen Bereichen ist entscheiden für die Wahl des Tiefpassfiltertyps. Es sollen nur Tschebyscheff-Filter sowie Butterworth-Filter betrachtet werden, da diese den einfachsten Aufbau haben. Der Butterworth-Filter besitzt im Durchlassbereich einen flachen Amplitudenverlauf. Darum wird dieser Filtertyp auch als maximal flach bezeichnet. Der Phasenverlauf ist in diesem Bereich nahezu linear. Diese Eigenschaft ist dabei bei System besonders wichtig, welche in Bezug auf ihre Phase besondere Genauigkeit erfordern, beispielsweise Winkelschätzung bei einem Radar. Im Sperrbereich erfolgt die Dämpfung monoton. Im Zeitbereich zeichnet sich der Butterworth-Filter durch starkes Überschwingen aus. Dieses beträgt bei einem Butterworth-Filter sechster Ordnung in etwa \SI{14}{\percent}(CITATION TIETZE, S.789).\\
Ein Tschebyscheff-Filter enthält im Durchlassbreich eine Welligkeit, sogenannte Ripple. Diese bewirken eine unterschiedliche betragsmäßige Übertragungsfunktion für verschiedene Frequenzen. Der Phasenverlauf kann nicht als linear angenommen werden. Demnach ist die Phasenänderung von der Frequenz abhängig. Im Sperrbereich nimmt die Dämpfung stärker zu als beim Butterworth-Filter. Ein Überschwingen im Zeitbereich findet nur in geringerem Ausmaß statt. In Abbildung XYZ sind beide Filtertypen in Bezugs auf den Amplitudengang sowie den Phasengang dargestellt. Die Grenze zwischen beiden Bereichern wird als Grenzfrequenz $f_{g}$ bezeichnet.
\section{Phasenregelschleife}
Eine Phasenregelschleife, auch PLL\footnote{ \textbf{P}hase \textbf{L}ocked \textbf{L}oop}, ist ein geschlossener Regelkreis der eine stabile Ausgangsfrequenz mit Hilfe einer Referenzfrequenz herstellt. Dabei wird eine PLL insbesondere bei der Taktgebung in Systemen verwendet. Bild \ref{fig:PLL_Blockerino} zeigt eine solche PLL als Blockschaltbild und wird weiterhin näher erläutert. Dabei wird ein Referenzfrequenzgeber, zwei Frequenzteiler, ein Phasendetektor, ein Schleifenfilter und ein Voltage Controlled Oscillator benötigt.
\begin{figure}[tbp]
  \centering
  \input{gfx/PLL_Block.tikz}
  \caption{Blockschaltbild Phasenregelkreis nach \cite{PLL_Perf}}
  \label{fig:PLL_Blockerino}
\end{figure}
Grundlage bei einer PLL ist stets eine stabile Referenz-Frequenz $f_{ref}$. Häufig wird diese durch einen Quarzoszillator erhalten. Der R-Teiler teilt die Frequenz in eine niedrigere, die Kompensationsfrequenz $f_{comp}$ mit 
\begin{align}
f_{comp} = \frac{f_{ref}}{R}.
\end{align} 
Diese Frequenz ist einer der Eingänge des Phasendetektors, PD. Der Phasendetektor gibt eine Spannung die dem Phasenfehler zwischen seinen beiden Eingängen, $f_{comp}$ und $f_{n}$, entspricht .  Dabei ist $f_{n}$ die durch den N-Teiler rückgeführte Ausgangsfrequenz. Die Charge Pump CP wandelt diese Spannung dabei in einen Strom $i_{CP}$ um. Der Phasendetektor hat dabei noch einen Proportionalität-Faktor $K_{\varphi}$. $K_{\varphi}$ ist dabei die Amplitude die die Charge-Pump verarbeiten kann. Wird $i_{CP}$ mit der Impedanz des Schleifenfilters $Z(s)$ multipliziert, was dem Ohmschen Gesetz entspricht, ergibt sich für den Voltage Controlled Oscillator VCO der Eingang 
\begin{align}
u_{VCO,in} = i_{CP}\cdot Z(s).
\end{align} 
Der Schleifenfilter ist dabei ein Tiefpass um lediglich die gewünschte Spannung an den VCO anzulegen. Der VCO ist ein Spannung-zu Frequenzwandler. Dabei hat der VCO einen Proportionalität-Faktor $K_{VCO}$. Demnach gilt für $f_{aus}$
\begin{align}
f_{aus} = \frac{R}{N} \cdot f_{Ref}.
\end{align} 
Durch ein hohes $R$ sowie ein kleines $N$ lassen sich somit Frequenzen in nahezu beliebiger Höhe erreichen, welches durch Referenzquarze nicht möglich ist, da diese häufig nur bis \SI{200}{\mega\hertz} verfügbar sind. Zusätzlich lässt sich dadurch eine niedrige Referenzfrequenz verwenden. Diese sind stabil und demnach ergibt sich geringeres Phasenrauschen am VCO.\\
Häufig wird der Begriff PLL für das gesamte Bild \ref{fig:PLL_Blockerino} verwendet, mit Ausnahme von Referenzquarz und VCO. Diese sind nach \cite[S.10]{PLL_Perf} schwer auf einen PLL Chip zu integrieren.\\
Entscheidende Parameter für die Leistung der PLL sind das Phasenrauschen und die Einrastzeit. Wesentlich für das Phasenrauschen ist dabei der N-Teiler Wert. Je kleiner dieser Wert desto kleiner das Phasenrauschen \cite[S.10]{PLL_Perf} . Die Einrastzeit ist die Dauer die die PLL für die Frequenzänderung benötigt, wenn der Wert des N-Teilers geändert wird.

%\chapter{Leitungstheorie und Fertigung}
In dieser Arbeit werden verschiedene Leiterplatten gefertigt. Eine solche Leiterplatte wird auch  PCB\footnote{\textbf{P}rinted \textbf{C}ircuit \textbf{B}oard} genannt. Diese Form der Leiterplattenherstellung ist bei Frequenzen von unter \SI{100}{\giga\hertz} üblich. Eine Leiterplatte ist aus verschiedenen Lagen aufgebaut. Entscheidend hierfür sind gewählte Dielektrika sowie die Materialdicken und die Lagenverteilung. Durch Ätzen, Bestrahlen und Kleben werden somit gewünschte Strukturen aufgebaut.  \\
Teil von PCBs sind immer planare Steifenleitungen. Die Signalführung auf planaren Steifenleitungen ist nur in seltenen Fällen exakt mathematisch berechenbar. Weiterhin wird ein Modell zur Berechnung einer Mikrostreifenleitung vorgestellt.
\section{Mikrostreifenleitung}
\begin{figure}[tp]
  \centering
  \includegraphics[scale=0.5]{gfx/MSL.png}
  \caption{Aufbau einer Mikrostreifenleitung\cite{MSL}}
  \label{fig:MSL}
\end{figure}
Eine Mikrostreifenleitung , dargestellt in Bild \ref{fig:MSL}, besteht aus einer leitenden Metallschicht auf der Oberseite eines Substrats. Unter diesem Substrat befindet sich ebenfalls eine leitende Metallschicht, die Massefläche\footnote{Auch Ground-Plane, die Bezugsfläche für das Potenzial}. Dabei sind die Substrat-Höhe $h$, die Metallbreite $W$, die Metallhöhe $t$ sowie die Dielektrizitätszahl $\epsilon_{r}$ des Substrats die charakteristischen Größen der Mikrostreifenleitung. Es muss für verschiedene Frequenzen von Signalen die Geometrie der Mikrostreifenleitung angepasst werden\cite{TransmissionLineDesignHandbook}. Da sich die elektromagnetische Welle überwiegend innerhalb des Substrats befindet, ist die effektive Dielektrizitätszahl $\epsilon_{r,eff}$ eine Möglichkeit die Feldverteilung mathematisch einfach darzustellen. Dabei ändert sich die ursprüngliche Geometrie der Mikrostreifenleitung mit Hilfe von $\epsilon_{r,eff}$ und der effektiven Leiterbahnbreite $W_{eff}$. Für $\epsilon_{r,eff}$ gilt dabei für $w/h < 1 $
\begin{align}
\epsilon_{r,eff} = \frac{\left( \epsilon_{r}+1 \right)}{2} + \frac{\left( \epsilon_{r}-1 \right)}{2} \cdot \left[ \left( 1 + \frac{10h}{W}\right)^{-0.5}+ 0.04 \left( 1.0 - \frac{w}{h}\right)^2\right]
\end{align}
sowie für $ w/h > 1$
\begin{align}
\epsilon_{r,eff} = \frac{\left( \epsilon_{r}+1 \right)}{2} + \frac{\left( \epsilon_{r}-1 \right)}{2} \cdot \left( 1 + \frac{10h}{W}\right)^{-0.5}.
\end{align}



\subsection{Differentielle Signalübertragung}
Eine differentielle Leitung besteht aus zwei Leiterbahnen. Auf jeder Leiterbahn ist ein Signal, welche zueinander komplementär sind, also $180^\circ $  Phasenverschiebung zueinander haben. Verglichen mit herkömmlichen Leitungen haben differentielle Leitungen einige Vorteile 
\begin{description}
\item Unterschiede im Bezugspotenzial zwischen Systemkomponenten haben kaum Einfluss auf Logiklevel.
\item Dämpfung auf einer Leiterbahn ist kompensierbar. Nach \cite{DifferentialSignalLine} lässt sich bis zu 24dB Dämpfung kompensieren.
\item Höhere Datenraten sind möglich. Dabei ist nach \cite{DifferentialSignalLine} bis zu  eine Datenrate von 10Gb/S möglich.
\end{description}
wobei stets mehr Leitungen nötig sind, um die gleiche Anzahl Signale zu führen, wenn pro Signal je eine Leitung verwendet wird. 
\begin{figure}[tbp]
  \centering
  \includegraphics[scale=1.5]{gfx/DiffSignalUebertragung.png}
  \caption{Einfluss einer Störung auf Differentielle Signalübertragung\cite{DiffSig}}
  \label{fig:DiffSig}
\end{figure}
In \ref{fig:DiffSig} ist eine differentielle Leitung dargestellt. Auf beide Signale wirkt dabei die gleiche Störung. Ist auf \ref{fig:DiffSig} die invertierte Leitung $s_{1}$ mit dem dazugehörigen Signal $-s(t)$ sowie die nicht invertierende Leitung $s_{2}$ mit $s(t)$ jeweils die Störung $e(t)$ addiert, ergibt sich durch die Differenz der am Operationsverstärker anliegenden Signale $f$ und $g$ mit 
\begin{align}
f(t) &= -s(t) + e(t) \\
g(t) &=  s(t) + e(t)
\end{align}
als Ausgang $s_{Res}(t)$ von
\begin{align}
s_{Res}(t) = g(t) - f(t)= 2\cdot s(t).
\end{align}
Dabei wird das SNR um 6dB verbessert.\\
Für die Leitungsanpassung der differentiellen Leitung muss die Impedanz der differentiellen Leitung $Z_{Diff}$ berechnet werden. Diese ergibt sich zu 
\begin{equation}
Z_{Diff} = 2Z_{0} \left( 1- 0.48e^{\left(-0.96\frac{s}{h} \right)} \right).
\end{equation}
Hier ist $s$ der Abstand zwischen den beiden einzelnen Leitungen. Häufig soll $Z_{Diff}$ = \SI{100}{\ohm} ergeben um optimal auf angepasst zu sein. Demnach kann dies bei entweder über $Z_{0}$ oder über $s$ eingestellt werden.Die jeweiligen Leitungslängen müssen gleich lang sein. Andernfalls ergibt sich ein Phasenunterschied der eine eindeutige Filterung verhindert, da dadurch gegebenenfalls Störungen zu unterschiedlichen Zeitpunkten anliegen. Je nach Frequenz des Signals verhindern bereits kleine Abweichungen die eindeutige Rekonstruktion des ursprünglichen Signals. Diese Abweichungen lassen sich leicht mit \ref{eq:fzulambda} berechnen. Ist die Toleranz der Leitungslänge etwa \SI{1}{\percent} so ist der zulässige Leitungslängenunterschied $\Delta l_{max}$ bei einer Frequenz von \SI{10}{\mega\hertz} $\Delta l_{max} = $ \SI{30}{\centi\meter}, bei einer Frequenz von \SI{2.4}{\giga\hertz} allerdings $\Delta l_{max} =$ \SI{1.25}{\milli\meter}. Der  Leitungslängenunterschied innerhalb einer Toleranz nimmt mit zunehmender Frequenz ab. Je nach Anwendung kann auch ein Bauteil-Takt die Toleranz vorgeben. Zusätzlich muss um eine Störung nur symmetrisch auf beide Leitungen zu erhalten der Leitungsabstand zueinander gering und konstant sein. Je nach Leitungsart ist dies allerdings auf Grund von Feldverteilungen nicht immer möglich.

\section{Lagenaufbau und Substrate}

Es werden Lagen entweder aus Kupfer oder aus einem Dielektrikum aufgebaut. Dabei dienen Kupferlagen insbesondere zur Abschirmung als Masse-Fläche oder Power/Ground-Lage. Zusätzlich sind Kupferlagen für die thermische Fluktuation, also zur Kühlung, geeignet. Besonders wichtig ist es, keine Inseln von Kupferlagen zu erzeugen. Diese Inseln haben kein definiertes Potenzial und verursachen potenziell Störungen. \\
Lagen aus einem Dielektrikum werden als Träger für Leiterbahnen verwendet. Dabei ist insbesondere die Dielektrizitätszahl entscheidend.\\
Es werden in dieser Arbeit drei verschiedene Substrate verwendet. Diese sind in Tabelle \ref{tab:FR4} dargestellt.
\begin{table}[tbp]
  \centering
  \begin{tabular}{lccr}
    Substratmaterial  & FR4 & Rogers4003 &Rogers3003\\
    \hline
    Dielektrizitätszahl $\epsilon_{r}$ & 5.0 & 3.38 & 3.0 \\
    Verlustfaktor $\tan \left( \delta \right)$ & 0.02 & 0.0021& 0.001 \\
    Oberflächenwiderstand & $\geq$\SI{10}{\kilo\ohm}& $\geq$\SI{4.2}{\giga\ohm} & $\geq$\SI{1}{\mega\ohm} \\
    Durchgangswiderstand & $\geq$\SI{10}{\kilo\ohm}&$\geq$\SI{17}{\giga\ohm}& $\geq$\SI{1}{\mega\ohm}\\
  \end{tabular}
  \caption{Materialparameter verwendeter Substrate}
  \label{tab:FR4}
\end{table}
FR4 ist ein gebräuchliches Substrat Material aus Epoxydharz-Glasfaser-Gewebe in der PCB-Herstellung. Die Eigenschaften von FR4 sind in Tabelle \ref{tab:FR4} dargestellt, wobei insbesondere der niedrige Preis und die einfache Verarbeitung, d.h. das Substrat ist fest und demnach nicht leicht verformbar, im Vergleich zu anderen Substraten von Vorteil ist.
\\
Ein Lagenwechsel, auch VIA, besteht aus einem PAD und einer zugehörigen Bohrung. Dabei sollen VIAs sowohl elektrisch als auch thermisch verbinden. Es gilt verschiedene Typen ,Bild \ref{fig:VIA}, von VIAs zu unterscheiden.
\begin{figure}[tbp]
  \centering
  \includegraphics[scale=0.5]{gfx/VIA.png}
  \caption{Verschiedene VIA\cite{VIA}}
  \label{fig:VIA}
\end{figure}
Dabei ist insbesondere das Through VIA weiterhin von Bedeutung. Durch dieses VIA lässt sich besonders gut thermische Leitfähigkeit sicherstellen und zusätzlich sind Through VIAs hervorragend als Messpunkte für die Inbetriebnahme von PCBs geeignet. \\
Wichtig sind die bei der Fertigung auftretenden Herstellungstoleranzen. Die Dicke einzelner Lagen kann sich um bis zu \SI{20}{\percent} verändern. Ebenfalls können sich einzelne Lagen zueinander verformen. Dabei sind besonders Innenlagen betroffen. Auch Bohr- oder Ätzungenauigkeiten sind zu berücksichtigen. Dabei sollten stets \SI{20}{\micro\meter} Toleranz bedacht werden. Demnach muss bei der Dimensionierung einer Komponente beachtet werden, dass die Größe um $\pm \SI{20}{\micro\meter}$ variiert. Dies tritt insbesondere bei Abständen von Komponenten zueinander auf. Jeder Abstand sollte demnach mindestens doppelt so groß sein wie die einfache Toleranz.
\subsection{Auslegung von Signalleitungen}
Signale oder Spannungen werden auf Kupfer geleitet. Dabei werden Spannungen eher über Ebenen geführt als über Leitungen mit konstante Breite. Dies resultiert aus der umgekehrten Proportionalität der Leitungsbreite zum Leitungswiderstand. Ein geringer Widerstand sorgt dabei für geringe Spannungsverluste auf dem PCB.\\
Signale werden entweder Single-Ended oder Differentiell geführt. Besonders digitale Signale sind dabei Single-Ended. Das eigene Zwischenfrequenzsignal wird auf dem PCB Differentiell geführt und erst mit der letzten Systemkomponente Single-Ended umgewandelt.\\
Die zu erwartende Bandbreite sind bis zu \SI{20}{\mega\hertz}. Bei diesen Frequenzen sind geringe Längenunterschiede in der differentielle Leitung tolerierbar, da bei angenommenen \SI{20}{\mega\hertz} ein Längenunterschied von \SI{8}{c\meter}, der gesamten Länge des PCB, nur ein relativer Fehler von 0.01875 oder von 10.6 Grad Phase auftritt. Da ein Fehler dieser Größenordnung bedeutet, dass eine Leitung über das gesamte PCB führt und die andere nicht existiert, kann der Fehler durch Längenunterschiede als klein genug angenommen werden um nicht weiter berücksichtigt zu werden.\\
Entsprechende Signalleitungen sind stets von Through-Vias umgeben. Diese sorgen für einheitliche Referenzpotenziale um die Leitungen herum. Zusätzlich sind Knicke von mehr als \SI{45}{\degree} zu vermeiden. Bei stärkerem Knick ändert sich abrupt die Impedanz der Leitung und verursacht demnach Reflexionen, die vermieden werden müssen. Theoretisch wären Leiterbahnbögen die beste Wahl, allerdings sind diese in der Herstellung deutlich aufwändiger im Vergleich zu \SI{45}{\degree} Winkeln und bringen dabei keine wesentliche Fortschritte in den betrachteten Frequenzen.\\
%\chapter{Gesamtkonzept}
\section{Aufbau}
Der Aufbau des Radarsystems ist Modular. Die einzelnen Bausteine sind dabei entsprechende Platinen mit definierter Funktion. Einzelne quadratischen Platinen haben eine Seitenlänge von \SI{8}{c\metre}.\\
 Verbunden sind die einzelnen Module durch \SI{1.27}{m\metre} SMC Connector der Firma ERNI. Diese ermöglichen Datenraten von bis zu \SI{3}{Gbit/s} und sind demnach für die geplanten Datenraten von maximal mehreren hundert MHz ausreichend. 
\begin{figure}[tbp]
  \centering
  \tikzset{%
	% Self defined bulding blocks. 
	% Nevertheless circutikz has implemented filters, couplers and other components since version 0.4, they are mostly implemented as bipoles.
	% The usage of bipoles: \draw (start) to[lowpass/amp/adc,....] (end).
	% The problem is, that if one wants to use arrows, the arrows in bipoles can not be sat manual (fixed in circuitikz source) AND THEY ARE NOT CONSISTENT
	% Also it's quite a mess, which component is a monopole, simple block, bipol, quad/triple etc
	% Following are a few examples on how to define your own blocks. 
	%
	% % % % % % % % % % % % % % % % % % % % % % % % % % % % % % % % % % % % % % % % % % % % % % % % % % % % % % % % % % % % % % % % % % % % %
	% % % % % % % % % % % % % % % % % % % % % % % % % % % % % % % % % % % % % % % % % % % % % % % % % % % % % % % % % % % % % % % % % % % % %
	% % % % % % % % % % % % % % % % % % % % % % % % % % % % % % % % % % % % % % % % % % % % % % % % % % % % % % % % % % % % % % % % % % % % %
	% % % % % % % % % % % % % % % % % % % % % % % % % % % % % % % % % % % % % % % % % % % % % % % % % % % % % % % % % % % % % % % % % % % % %
	%
	% Standard block definition, the width and height is adopted from the circutizk source code, so don't mind the strange values. Also the linewidth is set according to the circutrikz source code.
	block/.style    	= 	{draw, fill=white, thick, rectangle, minimum height = 0.98cm, minimum width = 0.98cm, node distance=2.5cm, line width=1.5pt},
	%
	% Standard circular block
	circleblock/.style	= 	{draw, fill=white, thick, circle, minimum width = 0.98cm,  line width=1.5pt, node distance=2.5cm},
	%
	% Label for circuitikz nodes, as they're reference is in the middle and not on the outer edge of the node....
	circuitikzlabel/.style	=	{label={[label, label distance=0.5cm]#1}},
	%
	%
	%
	% VCO/Oscillator 
	myVCO/.style			= 	{circleblock, path picture={%
		\draw[line width=0.75pt] 	($(path picture bounding box.west)+(0.09cm,0)$) sin ($(path picture bounding box.center)-(0.2cm,-0.2cm)$) cos  (path picture bounding box.center) sin ($(path picture bounding box.center)-(-0.2cm,0.2cm)$) cos ($(path picture bounding box.east)-(0.09cm,0)$);
		}
	},
	% Amplifier, as circuitikz does only provite amplifiers as 2-ports/bipoles
	myAMP/.style		= 	{block, node distance=2.5cm, path picture={%
		\draw[fill=white, line width=0.75pt] ($(path picture bounding box.center)+(0.7em,0)$) -- ($(path picture bounding box.center)-(0.7em,-0.7em)$) -- ($(path picture bounding box.center)-(0.7em,0.7em)$)  -- cycle;
		}
	},%%
	% Block	
	myBlock/.style    	= 	{draw, fill=white, thick, rectangle, minimum height = 0.98cm, minimum width = 0.98cm, node distance=2.5cm, line width=1.5pt},
	myBigBlock/.style    	= 	{draw, fill=white, thick, rectangle, minimum height = 0.98cm, minimum width = 2.94cm, node distance=2.5cm, line width=1.5pt},	
	% Same for ADC
	myADC/.style 	=	{block, path picture={%
		\draw[line width=0.75pt] 	(path picture bounding box.south west) -- (path picture bounding box.north east);
		\node[] at ($(path picture bounding box.center)+(-.5em,.5em)$) () {D};
		\node[] at ($(path picture bounding box.center)+(.5em,-.5em)$) () {A};
		} 
	},
	% Same for filters
	myLP/.style	=	{block, path picture={%
		%Sine-Waves
		\draw[line width=.75pt] 	($(path picture bounding box.west)+(0.3em,0)$) sin ($(path picture bounding box.center)-(0.50em,-0.3em)$) cos  (path picture bounding box.center) sin ($(path picture bounding box.center)-(-0.50em,0.3em)$) cos ($(path picture bounding box.east)-(0.3em,0)$);
		\draw[line width=0.75pt] 	($(path picture bounding box.west)+(0.3em,-0.65em)$) sin ($(path picture bounding box.center)-(0.50em,0.35em)$) cos  ( $(path picture bounding box.center)-(0,0.65em)$) sin ($(path picture bounding box.center)-(-0.50em,0.95em)$) cos ($(path picture bounding box.east)-(0.3em,0.65em)$);
		\draw[line width=0.75pt] 	($(path picture bounding box.west)+(0.3em,0.65em)$) sin ($(path picture bounding box.center)-(0.50em,-0.95em)$) cos  ( $(path picture bounding box.center)+(0,0.65em)$) sin ($(path picture bounding box.center)-(-0.50em,-0.35em)$) cos ($(path picture bounding box.east)-(0.3em,-0.65em)$);
		% Cancelation
		\draw[line width=0.75pt] 	($(path picture bounding box.center)-(0.2em,0.2em)$) -- (path picture bounding box.center) -- ($(path picture bounding box.center)+(0.2em,0.2em)$) ;
		\draw[line width=0.75pt] 	($(path picture bounding box.center)-(0.2em,-0.45em)$) -- ($(path picture bounding box.center)+(0,0.65em)$) -- ($(path picture bounding box.center)+(0.2em,0.85em)$) ;
		}
	},
}
\begin{tikzpicture}[line width=0.7pt,>=latex,node distance=2.5cm,scale = 0.7]

	%Box
	\draw[dashed,blue,-] (-1,1.5) rectangle (13,15.5);
	\draw(6,15.5)
		node[label={above:Radarsystem}](){};
	%ADC
	\draw (-4,4.5)
		node[myADC, label={below:Analog-Digital-Wandler}] (ReFF) {};
	\draw[->](0,4.5) -- (ReFF.east); 
	
	%RefTakt
	\draw (-4,6.5)
		node[myVCO, label={above:Referenztakt}] (ReF) {};
	\draw[<-](0,6.5) -- (ReF.east); 
	
	%Spannungsversorgung
	\draw[->](-3,2.5) -- (0,2.5);
	\draw (-4,3) 
		node[label={below:$\text{V}_{\text{Supply}}$}](abc){};
		
	%Radar-Board
	\draw[-] (6,8)--(12,8);
	\draw[-](12,8)--(12,9);
	\draw[-](12,9)--(6,9);
	\draw[-](6,9)--(6,8);		
	\draw(9,9)
		node[label ={below:Radar-Platine}] (Ref133337) {};
			
			
	%PLL-Board
	\draw[-](0,7)--(0,6);
	\draw[-](0,6)--(12,6);
	\draw[-](12,6)--(12,7);
	\draw[-](12,7)--(0,7);	
	\draw(6,7)
		node[label ={below:Phasenregelkreis-Platine}] (Ref13337) {};
	
	%IF-Board
	\draw[-](0,5)--(0,4);
	\draw[-](0,4)--(12,4);
	\draw[-](12,4)--(12,5);
	\draw[-](12,5)--(0,5);
	\draw(6,5)
		node[label ={below:Zwischenfrequenz-Platine}] (Ref133337) {};
	
	%Power-Board
	\draw[-](0,3)--(0,2);
	\draw[-](0,2)--(12,2);
	\draw[-](12,2)--(12,3);
	\draw[-](12,3)--(0,3);
	\draw(6,3)
		node[label ={below:Power-Platine}] (Ref1337) {};

	%Linse
	\draw[-](0,12)--(6,12);
	\draw[-](6,12)--(6,13);
	\draw[-](0,13)--(0,12);
	\draw[-](0,13)--(1,13);
	\draw[-](6,13)--(5,13); 		
	\draw(3,13)
		node[label ={below:Dieelektrische Linse}] (Ref1333337) {};	
	\draw (5,13) arc(0:180:2cm);
	
	%Antenna
	\draw(3,8.5)
		node[antenna](){};
	\draw[dashed,-](3,8.5)--node[above]{Optional}(6,8.5);
	

\end{tikzpicture}

  \caption{\SI{122}{G\hertz} Radarsystem - schematischer Aufbau}
  \label{fig:Gesamtkonzept_Grob}
\end{figure} 
Der maximale Storm pro Pin sind \SI{1.7}{\ampere}. Die Einfügedämpfung ist laut Datenblatt in einem Bereich bis \SI{4}{\giga\hertz} stets unter \SI{0.4}{dB}. Übersprechen ist  bei nebeneinander liegenden differentiellen Signalen im Bereich bis \SI{4}{\giga\hertz} stets unter \SI{-15}{dB}. Bei Frequenzen im MHz-Bereich ist Übersprechen stets unter \SI{-30}{dB} und somit hat die Pinbelegung praktisch keinen Einfluss auf die Signalqualität. 
Es wurden pro benötigte Platinenseite 4 Pinleisten verwendet. Dabei je zwei Leisten mit 26 Pins und zwei Leisten mit 12 Pins. Falls sowohl die Unterseite als auch die Oberseite des Boards beschaltet wird sind auf beiden Seiten 4 Pinleisten vorzufinden. Dementsprechend benötigen die unterste sowie die oberste Platine nur 4 Pinleisten um die Funktionalität mit den restlichen \SI{8}{c\metre} Platinen gewährleisten.\\
Das Radarsystem ist nicht in sich geschlossen. Verbindungen aus dem System werden über SMA Verbinder der Firma Molex hergestellt. Diese sind bei RF-Anwendungen üblich. \\
In Bild \ref{fig:Gesamtkonzept_Grob} ist der schematische Aufbau dargestellt. Dabei bildet der Mikrocontroller das unterste Modul. Auf dem Mikrocontroller ist die Versorgungsplatine aufgesteckt. Diese stellt die für das System benötigten Ströme und Spannungen bereit. Darüber ist die Zwischenfrequenzplatine, welche die Signalkonditionierung vor dem ADC vornimmt. Darüber ist die Phasenregelkreisplatine. Diese stellt die für den Radarchip benötigten Signale bereit. Darauf befindet sich die kleinere Radarplatine. Diese ist von den Abmessungen kleiner und wird separat vorgestellt. Optional ist eine dieeletrische Linse zu platzieren. \\
Die Signalkonditionierung verwendet in erster Stufe ein ADF4937-2 Operationsverstärker ausgelegt auf \SI{10}{dB} Gain. Die zweite Stufe ist ein programmierbarer HMC960 Verstärker der zwischen \SI{0}{dB} und \SI{40}{dB} Gain in \SI{0.5}{dB} Schritten erzeugt. Die dritte Stufe ist ein HMC900 ein programmierbarer Tiefpass mit \SI{0}{dB} oder \SI{10}{dB} Gain. Als vierte und letzte Systemstufe ist ein AD8034 Operationsverstärker. Dieser konvertiert die bis dahin differentiell geführten Signale in einfache, auf einer Datenbahn geführten Signale. Der AD8034 ist mit \SI{10}{dB} auslegt. Sämtliche Komponenten sind von der Firma Analog Devices.\\
Der Phasenregelkreis wird durch eine Kombination des ADF4159 der Firma Analog Devices zur Frequenzgenerierung sowie einem LTC6954 der Firma Linear Technology als Takt-Teiler in Kombination mit dem spannungsgesteuerten
Oszillator  auf dem Radarchip hergestellt. Dabei ist der ADF4159 zur Rampengenerierung notwendig. Optional soll der LTC6954 entweder über ein Referenzquarz oder über einen externen Takt angesteuert werden.
\section{Radarchip}
Der in dieser Arbeit verwendete Radarchip(TRX 120 06 von Silicon-Radar) ist ein Sende-Empfangssystem für das ISM-Band(\SI{122}{G\hertz}.  Low-Noise-Ampflifier, Quadratur-Mischstufe, Mehrphasenfilter, VCO sowie 1/32 Frequenzteiler sind bereits Teil des Systems. 
\subsection{Funktionsweise Radarchip}
Das Blockschaltbild ist in Bild \ref{fig:Radarchip} dargestellt. Der spannungsgesteuerte Oszillator wird durch Grob-und Feinabstimmung eingestellt. Dabei mit einer 3 Bit großen Grobabstimmung $\text{V}_{\text{grob}}$ sowie mit einer 1 Bit großen Feinabstimmung $\text{V}_{\text{fein}}$ die gewünschte Frequenz eingestellt werden. Durch entsprechende Kombination der Eingänge können sehr große Bandbreiten von bis zu \SI{7}{G\hertz} erreicht werden. Zusätzlich kann der Sendesignalpegel mit Hilfe von zwei Leistungsdetektoren zwischen dem Verstärker und der Sendeantenne gemessen werden. Das Empfangssignal wird zunächst durch einen LNA verstärkt und anschließend mit Hilfe der Quadratur-Mischstufe in das Basisband konvertiert. Diese beiden Signale, $\text{IF}_{\text{I}}$ und $\text{IF}_{\text{Q}}$, werden dabei jeweils differentiell geführt. Ein Teil der Oszillator-Frequenz wird durch den 1/32 Frequenzteiler differentiell ausgegeben. Dieser Ausgang kann, falls die Verwendung ein FMCW-Radar sein soll, in einen Phasenregelkreis eingegeben werden und damit die Funktionalität sicherstellen.
\begin{figure}[tbp]
  \centering
  \tikzset{%
	% Self defined bulding blocks. 
	% Nevertheless circutikz has implemented filters, couplers and other components since version 0.4, they are mostly implemented as bipoles.
	% The usage of bipoles: \draw (start) to[lowpass/amp/adc,....] (end).
	% The problem is, that if one wants to use arrows, the arrows in bipoles can not be sat manual (fixed in circuitikz source) AND THEY ARE NOT CONSISTENT
	% Also it's quite a mess, which component is a monopole, simple block, bipol, quad/triple etc
	% Following are a few examples on how to define your own blocks. 
	%
	% % % % % % % % % % % % % % % % % % % % % % % % % % % % % % % % % % % % % % % % % % % % % % % % % % % % % % % % % % % % % % % % % % % % %
	% % % % % % % % % % % % % % % % % % % % % % % % % % % % % % % % % % % % % % % % % % % % % % % % % % % % % % % % % % % % % % % % % % % % %
	% % % % % % % % % % % % % % % % % % % % % % % % % % % % % % % % % % % % % % % % % % % % % % % % % % % % % % % % % % % % % % % % % % % % %
	% % % % % % % % % % % % % % % % % % % % % % % % % % % % % % % % % % % % % % % % % % % % % % % % % % % % % % % % % % % % % % % % % % % % %
	%
	% Standard block definition, the width and height is adopted from the circutizk source code, so don't mind the strange values. Also the linewidth is set according to the circutrikz source code.
	block/.style    	= 	{draw, fill=white, thick, rectangle, minimum height = 0.98cm, minimum width = 0.98cm, node distance=2.5cm, line width=1.5pt},
	%
	% Standard circular block
	circleblock/.style	= 	{draw, fill=white, thick, circle, minimum width = 0.98cm,  line width=1.5pt, node distance=2.5cm},
	%
	% Label for circuitikz nodes, as they're reference is in the middle and not on the outer edge of the node....
	circuitikzlabel/.style	=	{label={[label, label distance=0.5cm]#1}},
	%
	%
	%
	% VCO/Oscillator 
	myVCO/.style			= 	{circleblock, path picture={%
		\draw[line width=0.75pt] 	($(path picture bounding box.west)+(0.09cm,0)$) sin ($(path picture bounding box.center)-(0.2cm,-0.2cm)$) cos  (path picture bounding box.center) sin ($(path picture bounding box.center)-(-0.2cm,0.2cm)$) cos ($(path picture bounding box.east)-(0.09cm,0)$);
		}
	},
	% Amplifier, as circuitikz does only provite amplifiers as 2-ports/bipoles
	myAMP/.style		= 	{block, node distance=2.5cm, path picture={%
		\draw[fill=white, line width=0.75pt] ($(path picture bounding box.center)+(0.7em,0)$) -- ($(path picture bounding box.center)-(0.7em,-0.7em)$) -- ($(path picture bounding box.center)-(0.7em,0.7em)$)  -- cycle;
		}
	},%%
	% Block	
	myBlock/.style    	= 	{draw, fill=white, thick, rectangle, minimum height = 0.98cm, minimum width = 0.98cm, node distance=2.5cm, line width=1.5pt},
	myBigBlock/.style    	= 	{draw, fill=white, thick, rectangle, minimum height = 0.98cm, minimum width = 2.94cm, node distance=2.5cm, line width=1.5pt},	
	% Same for ADC
	myADC/.style 	=	{block, path picture={%
		\draw[line width=0.75pt] 	(path picture bounding box.south west) -- (path picture bounding box.north east);
		\node[] at ($(path picture bounding box.center)+(-.5em,.5em)$) () {D};
		\node[] at ($(path picture bounding box.center)+(.5em,-.5em)$) () {A};
		} 
	},
	% Same for filters
	myLP/.style	=	{block, path picture={%
		%Sine-Waves
		\draw[line width=.75pt] 	($(path picture bounding box.west)+(0.3em,0)$) sin ($(path picture bounding box.center)-(0.50em,-0.3em)$) cos  (path picture bounding box.center) sin ($(path picture bounding box.center)-(-0.50em,0.3em)$) cos ($(path picture bounding box.east)-(0.3em,0)$);
		\draw[line width=0.75pt] 	($(path picture bounding box.west)+(0.3em,-0.65em)$) sin ($(path picture bounding box.center)-(0.50em,0.35em)$) cos  ( $(path picture bounding box.center)-(0,0.65em)$) sin ($(path picture bounding box.center)-(-0.50em,0.95em)$) cos ($(path picture bounding box.east)-(0.3em,0.65em)$);
		\draw[line width=0.75pt] 	($(path picture bounding box.west)+(0.3em,0.65em)$) sin ($(path picture bounding box.center)-(0.50em,-0.95em)$) cos  ( $(path picture bounding box.center)+(0,0.65em)$) sin ($(path picture bounding box.center)-(-0.50em,-0.35em)$) cos ($(path picture bounding box.east)-(0.3em,-0.65em)$);
		% Cancelation
		\draw[line width=0.75pt] 	($(path picture bounding box.center)-(0.2em,0.2em)$) -- (path picture bounding box.center) -- ($(path picture bounding box.center)+(0.2em,0.2em)$) ;
		\draw[line width=0.75pt] 	($(path picture bounding box.center)-(0.2em,-0.45em)$) -- ($(path picture bounding box.center)+(0,0.65em)$) -- ($(path picture bounding box.center)+(0.2em,0.85em)$) ;
		}
	},
}
\begin{tikzpicture}[line width=0.7pt,>=latex,node distance=2.5cm,scale = 0.8]
	% First: All building blocks are placed relative to the first component
	\draw (0,0)
		node[myVCO, label={below: Oszillator}] (Ref) {};
	\draw (0,2)
		node[myBlock] (Ref1) {1/32};
	\draw (2,0)
		node[myAMP] (Ref2) {};
	\draw (4,2)
		node[myAMP] (Ref3) {};
	\draw (4,0)
		node[myAMP] (Ref4) {};
	\draw (6,-2)
		node[myBlock] (Ref5) {90$^\circ$};
	\draw (8,-4)
		node[mixer] (Ref6) {};
	\draw (8,-6)
		node[mixer] (Ref7) {};
	\draw (10,-4)
		node[myAMP, rotate = 180, label={above:LNA}] (Ref8) {};
	\draw (11,-4)
		node[rxantenna, label={below right:RX}] (Ref10) {};
	\draw (11,2)
		node[txantenna, label={below right:TX}] (Ref9) {};
		
	%Connections
	\draw[->](Ref)--(Ref1);
	\draw[->](Ref)--(Ref2);
	\draw[-](Ref2)--(3,0);
	\draw[->](3,0)|-(Ref3);
	\draw[->](3,0)--(Ref4);
	\draw[->](Ref4)-|(Ref5);
	\draw[->](Ref4)-|(Ref6.north);
	\draw[-](Ref5)|- (8,-5);
	\draw[->](8,-5)--(Ref7.north);
	\draw[-](Ref8.east)--(9,-4);
	\draw[->](9,-4)--(Ref6.east);
	\draw[->](9,-4)|-(Ref7.east);
	\draw[-](Ref10)--(Ref8.west);
	\draw[-](Ref6.west)--(0,-4);
	\draw[-](Ref7.west)--(0,-6);
	\draw[->](0,-4)--node[above left]{$\text{IF}_{\text{I}}$}node[below left]{(differentiell)}(-2,-4);
	\draw[->](0,-6)--node[above left]{$\text{IF}_{\text{Q}}$}node[below left]{(differentiell)}(-2,-6);
	\draw[-](Ref9)--(Ref3);
	
	
	%Div_o
	\draw[->](Ref1)--node[above left]{$\text{div}_{\text{out}}$}node[below left]{(differentiell)}(-2,2);
	%Oszi In
	\draw[->](-2,0)--node[above left]{$\text{V}_{\text{fein}}$}node[below left]{$\text{V}_{\text{grob}}$, 3Bit}(Ref);
		
	
	%Blue Box
	\draw[dashed,blue,-](-1,8)--(-1,-8);
	\draw[dashed,blue,-](-1,-8)--(15,-8);
	\draw[dashed,blue,-](-1,8)--(15,8);
	\draw[dashed,blue,-](15,8)--(15,-8);	
	
	%Control Unit
	\draw[-](0,7)--(0,5);
	\draw[-](0,5)--(9,5);
	\draw[-](9,5)--(9,7);
	\draw[-](9,7)--(0,7);
	%Output Control Unit
	\draw[->](0,6)--node[above left]{deto}node[below left]{(differentiell)}(-2,6);	
	%Input Control Unit	
	\draw[-](1,7)--node[above left]{s3}(1,9);
	\draw[-](2,7)--node[above left]{s2}(2,9);
	\draw[-](3,7)--node[above left]{s1}(3,9);
	\draw[-](4,7)--node[above left]{s0}(4,9);
	\draw[-](6,7)--node[above left]{$\text{div}_{\text{en}}$}(6,9);
	\draw[-](8,7)--node[above left]{$\text{buf}_{\text{en}}$}(8,9);
	%Text Control Unit
	\draw(4.5,6.5)
		node[label ={below:Control Unit}] (Ref1337) {};


\end{tikzpicture}

  \caption{Blockschaltbild \SI{122}{G\hertz}Radarchip(QUELLE)}
  \label{fig:Radarchip}
\end{figure}
\subsection{Anwendung}
Der Radarchip kann sowohl als FMCW-Radar auch als CW-Radar ,über die Beschaltung des Oszillators, betrieben des Oszillator werden. Bei konstanter Oszillator-Frequenz entspricht dies dem Betrieb als CW-Radar, andernfalls wird der Radarchip als FMCW-Radar betrieben. \\
Der TXR kann für Reichweiten von bis zu \SI{10}{m} verwendet werden. Mit Hilfe dieelektrischer Linsen kann die Reichweite zusätzlich erhöht werden. Als FMCW-Radar sind Bandbreiten bis \SI{7}{G\hertz} möglich, wobei der Chip auf die Verwendung im ISM Band von \SI{122}{G\hertz} bis \SI{123}{G\hertz} ausgelegt ist.
\subsection{Versorgung und Beschaltung}
Der Radarchip wird auf eine eigene Platine aufgesetzt, bildet dadurch die oberste Komponente des Radarsystems, und durch diese beschaltet. Die Verbindung zu dem darunterliegenden Board wird über zwei zweireihige \SI{2.0}{m\metre} Pinleisten mit jeweils insgesamt 20 Pins hergestellt.\\
Beide Zwischenfrequenzsignale $\text{IF}_{\text{I}}$ und $\text{IF}_{\text{Q}}$, der Teiler Ausgang $\text{div}_{\text{out}}$ sowie der Multiplexer Ausgang deto werden differentiell geführt. $\text{div}_{\text{out}}$ wird über MML ausgeführt. Deto wird durch die vier Kontrollsignale s0, s1,  s2, und s3 gesteuert. $\text{div}_{\text{en}}$ und $\text{buf}_{\text{en}}$ schalten dabei den Frequenzteiler und die Buffer ein. Die Kontrollspannung am Oszillator wird durch ein Widerstandsnetzwerk eingestellt. Einzelne Widerstände können dabei je nach Bedarf als Bestückoption offen gelassen werden.\\
Das Radarsystem kann mit der im Chip integrierten monostatischen Antenne beschaltet werden. Es ist auch möglich, extern Antennen anzuschließen. Dieser Fall ist in Bild \ref{fig:Gesamtkonzept_Grob} dargestellt. \\
Wesentlich für die Funktionalität des Radarsystems ist der Referenztakt. Damit die Kombination aus Radar-Platine und Phasenregelkreis funktioniert, ist geringes Phasenrauschen erforderlich, da sich beide Komponenten auf den Referenztakt beziehen. Falls dort Phasenfehler auftreten, ist Kohärenz nicht gegeben.
\section{Mikrocontroller}
Die Radarsteuerung wird über einen Mikrocontroller, ein STM32f303, vorgenommen. Dieser 32 Bit Mikrocontroller kann bis zu \SI{84}{M\hertz} takten . Die SPI-Interfaces können Datenraten von bis zu \SI{42}{Mbit/s} verarbeiten.
%\chapter{Zwischenfrequenzplatine}
In diesem Kapitel wird die Signalkonditionierung beschrieben. Dabei wird die Signalkonditionierung auf einer seperaten Platine, die Zwischenfrequenzplatine, vorgenommen. Eingangssignal werden zwei Kanäle, I und Q, in jeweils differentieller Form sein. Diese werden direkt vom Radarchip auf das Board geführt. Ziel ist es, als Ausgangssignal ein für die ADC Wandlung konditioniertes Signal zu erhalten. Ebenfalls soll die Zwischenfrequenzplatine unabhängig vom verwendeten Radarchip sein. 
\section{Anforderung und Verwendung}
Ein wesentliches Kriterium ist die erwartete Bandbreite. Diese ist abhängig von Radarchip und den darin verwendeten Mischern. Für den primär verwendeten \SI{122}{G\hertz} Chip ist eine Zwischenfrequenzbandbreite von 0-\SI{20}{M\hertz} angegeben. Dabei ist die Ausgangsimpedanz ausgehend vom Radarchip \SI{500}{\ohm}. Weiterhin muss der Eingangspegel bestimmt werden. Minimal liegt bei nach (LINK BUDGET REFERENZ) bei \SI{-44}{dBm}.
\section{Auswahl und Anordnung der Komponenten}
In Bild \ref{fig:IFBoard_Block} ist das Blockdiagramm der Zwischenfrequenzplatine dargestellt. Der Eingang links im Bild ist das vom Radarchip ausgehende Signal. Der Ausgang rechts ist das an den ADC weitergegebene, signalkonditionierte Signal. Die ersten drei Systemkomponenten sind differentiell zu differentiell. Demnach besitzen sie jeweils pro Kanal zwei Eingänge und zwei Ausgänge. Die letzte Systemstufe hat einen differentiellen Eingang und einen einfachen Ausgang.\\
Als erste Komponente wurde ein ADA4937-2 von Analog Devices gewählt. Dieser differentiell zu differentiell Verstärker wird mit \SI{5}{\volt} betrieben. Der ADA4837-2 hat eine hohe differentielle Eingangsimpedanz von \SI{6}{M\ohm} sowie eine auf den Gleichanteil bezogene Eingangsimpedanz von \SI{3}{M\ohm}. Dies sorgt für Stabilität bei den ausgehenden Radarchipsignalen. Zusätzlich ist die Gleichtaktunterdrückung \SI{-80}{dB} wodurch Rauschen auf dem Gleichtakt stets unterdrückt wird. Die \SI{-3}{dB} Bandbreite liegt bei \SI{1.9}{G\hertz} und ist somit mehr als ausreichend. Der Gain kann über ein Widerstandsnetzwerk von vier Widerständen eingestellt werden. \\
\begin{figure}[tbp]
  \centering
  \input{gfx/IfBoard_Block.tikz}
  \caption{Blockdiagramm Ziwschenfrequenzplatine}
  \label{fig:IFBoard_Block}
\end{figure} 
Die zweite Komponente ist ein HMC960 von Analog Devices. Dieser in Kombination mit der dritten Komponente, einem HMC900, ist die eigentliche Signalkonditionierung bestehend aus einem Variablen-Gain-Verstärker (HMC960) sowie einem Variablen-Tiefpass (HMC900). Der HMC900 hat dabei noch eine optionale \SI{10}{dB} Gain Stufe.\\
Die differentiell zu einfach Konvertierung wird von einem AD8034 vorgenommen. Dieser ist die vierte und letzte Systemkomponente.
\section{Filterentwurf}
Um Rauschen auf den Gleichanteil zu unterdrücken werden zwischen den einzelnen Systemkomponenten Filter eingefügt, die bei Bedarf bestückt werden können. Zusätzlich soll das eigentliche Signal nicht gedämpft werden. Die Bandbreite von \SI{20}{M\hertz} muss demnach eingehalten werden.
\begin{figure}[tbp]
  \centering
  \tikzset{%
	% Self defined bulding blocks. 
	% Nevertheless circutikz has implemented filters, couplers and other components since version 0.4, they are mostly implemented as bipoles.
	% The usage of bipoles: \draw (start) to[lowpass/amp/adc,....] (end).
	% The problem is, that if one wants to use arrows, the arrows in bipoles can not be sat manual (fixed in circuitikz source) AND THEY ARE NOT CONSISTENT
	% Also it's quite a mess, which component is a monopole, simple block, bipol, quad/triple etc
	% Following are a few examples on how to define your own blocks. 
	%
	% % % % % % % % % % % % % % % % % % % % % % % % % % % % % % % % % % % % % % % % % % % % % % % % % % % % % % % % % % % % % % % % % % % % %
	% % % % % % % % % % % % % % % % % % % % % % % % % % % % % % % % % % % % % % % % % % % % % % % % % % % % % % % % % % % % % % % % % % % % %
	% % % % % % % % % % % % % % % % % % % % % % % % % % % % % % % % % % % % % % % % % % % % % % % % % % % % % % % % % % % % % % % % % % % % %
	% % % % % % % % % % % % % % % % % % % % % % % % % % % % % % % % % % % % % % % % % % % % % % % % % % % % % % % % % % % % % % % % % % % % %
	%
	% Standard block definition, the width and height is adopted from the circutizk source code, so don't mind the strange values. Also the linewidth is set according to the circutrikz source code.
	block/.style    	= 	{draw, fill=white, thick, rectangle, minimum height = 0.98cm, minimum width = 0.98cm, node distance=2.5cm, line width=1.5pt},
	%
	% Standard circular block
	circleblock/.style	= 	{draw, fill=white, thick, circle, minimum width = 0.98cm,  line width=1.5pt, node distance=2.5cm},
	%
	% Label for circuitikz nodes, as they're reference is in the middle and not on the outer edge of the node....
	circuitikzlabel/.style	=	{label={[label, label distance=0.5cm]#1}},
	%
	%
	%
	% VCO/Oscillator 
	myVCO/.style			= 	{circleblock, path picture={%
		\draw[line width=0.75pt] 	($(path picture bounding box.west)+(0.09cm,0)$) sin ($(path picture bounding box.center)-(0.2cm,-0.2cm)$) cos  (path picture bounding box.center) sin ($(path picture bounding box.center)-(-0.2cm,0.2cm)$) cos ($(path picture bounding box.east)-(0.09cm,0)$);
		}
	},
	% Amplifier, as circuitikz does only provite amplifiers as 2-ports/bipoles
	myAMP/.style		= 	{block, node distance=2.5cm, path picture={%
		\draw[fill=white, line width=0.75pt] ($(path picture bounding box.center)+(0.7em,0)$) -- ($(path picture bounding box.center)-(0.7em,-0.7em)$) -- ($(path picture bounding box.center)-(0.7em,0.7em)$)  -- cycle;
		}
	},%%
	% Block	
	myBlock/.style    	= 	{draw, fill=white, thick, rectangle, minimum height = 0.98cm, minimum width = 0.98cm, node distance=2.5cm, line width=1.5pt},
	myBigBlock/.style    	= 	{draw, fill=white, thick, rectangle, minimum height = 0.98cm, minimum width = 2.94cm, node distance=2.5cm, line width=1.5pt},	
	% Same for ADC
	myADC/.style 	=	{block, path picture={%
		\draw[line width=0.75pt] 	(path picture bounding box.south west) -- (path picture bounding box.north east);
		\node[] at ($(path picture bounding box.center)+(-.5em,.5em)$) () {D};
		\node[] at ($(path picture bounding box.center)+(.5em,-.5em)$) () {A};
		} 
	},
	% Same for filters
	myLP/.style	=	{block, path picture={%
		%Sine-Waves
		\draw[line width=.75pt] 	($(path picture bounding box.west)+(0.3em,0)$) sin ($(path picture bounding box.center)-(0.50em,-0.3em)$) cos  (path picture bounding box.center) sin ($(path picture bounding box.center)-(-0.50em,0.3em)$) cos ($(path picture bounding box.east)-(0.3em,0)$);
		\draw[line width=0.75pt] 	($(path picture bounding box.west)+(0.3em,-0.65em)$) sin ($(path picture bounding box.center)-(0.50em,0.35em)$) cos  ( $(path picture bounding box.center)-(0,0.65em)$) sin ($(path picture bounding box.center)-(-0.50em,0.95em)$) cos ($(path picture bounding box.east)-(0.3em,0.65em)$);
		\draw[line width=0.75pt] 	($(path picture bounding box.west)+(0.3em,0.65em)$) sin ($(path picture bounding box.center)-(0.50em,-0.95em)$) cos  ( $(path picture bounding box.center)+(0,0.65em)$) sin ($(path picture bounding box.center)-(-0.50em,-0.35em)$) cos ($(path picture bounding box.east)-(0.3em,-0.65em)$);
		% Cancelation
		\draw[line width=0.75pt] 	($(path picture bounding box.center)-(0.2em,0.2em)$) -- (path picture bounding box.center) -- ($(path picture bounding box.center)+(0.2em,0.2em)$) ;
		\draw[line width=0.75pt] 	($(path picture bounding box.center)-(0.2em,-0.45em)$) -- ($(path picture bounding box.center)+(0,0.65em)$) -- ($(path picture bounding box.center)+(0.2em,0.85em)$) ;
		}
	},
}
\begin{tikzpicture}[line width=0.7pt,>=latex,node distance=2.5cm,scale = 0.8]
	% First: All building blocks are placed relative to the first component
	
	\draw (0,0) to[R=$R_1$] (2,0);	
	\draw (0,-2) to[R=$R_2$] (2,-2);	
	\draw (3,2) to[C=$C_1$] (3,0);	
	\draw (3,-2) to[C=$C_2$] (3,-4);	
	\draw (4,0) to[C=$C_3$] (4,-2);

	\draw[o-](-2,0)--(0,0);
	\draw[o-](-2,-2)--(0,-2);
	
	\draw[-o] (2,0)--(6,0);
	\draw[-o] (2,-2)--(6,-2);
	
	\draw(3,2)node[ground,rotate = 180](){};
	\draw(3,-4)node[ground](){};
\end{tikzpicture}

  \caption{Blockdiagramm CM-Filter}
  \label{fig:CM_Filter_Block}
\end{figure} 
\section{Schaltplan-Erstellung und PCB-Design}
\begin{figure}[tbp]
  \centering
  \input{gfx/CM_Filter_Simulation.tikz}
  \caption{Blockdiagramm CM-Filter}
  \label{fig:CM_Filter_Block}
\end{figure} 
\section{Softwarekonzept}






% Testinhalt laden, dafür muss Paket 'blindtext' in 'header.tex' geladen werden
%\blinddocument % kleines Dokument
%\Blinddocument % großes Dokument


%% Anhang einleiten
%\appendix


%% Verzeichnisse
%\chapter*{Symbole}
%\addcontentsline{toc}{chapter}{Symbole}
%\glsaddall[types={symblist}]
%\glsaddall[types={konvlist}]
%\printglossary[type=konvlist, style=long]
%\printglossary[type=symblist,style=long]

%\listoffigures
%\listoftables

\bibliography{literatur}
% \printbibliography

\ifcase \LANGUAGE
    % \bibliographystyle{unsrt}                   % do not use 'unsrt' anymore, due to compatibility issues
    \bibliographystyle{bibtex/IEEEtran}         % use the IEEE default
    % \bibliographystyle{bibtex/IEEEtranSA}       % a IEEE style but with alphanumeric behaviour
\else
    \bibliographystyle{bibtex/myIEEEtran-de}    % use the IEEE default, german adaptions by Tobias Chaloun
    % \bibliographystyle{alphadin}
    % \bibliographystyle{unsrtdin}
    % \bibliographystyle{gerunsrt}                % do not use 'gerunsrt' anymore, due to compatibility issues
\fi

\end{document}
